%% ----------------------------------------------------------------
%% Tesis.tex -- Fichero Principal (el que se compila con Latex)
%% ---------------------------------------------------------------- 

%% Esta plantilla está basada en la Tesis de Grado escrita por Sunil Patel,
% (quién la basó en la plantilla ecsthesis) bajo la Licencia Pública del Proyecto LaTeX
% que puede encontrarse acá: http://latex-project.org/lppl/
% con la esperanza de que sea más fácil de usar y escalar a sus necesidades
% por Simon Ternsjö en 2013-10
% Traducida al español para su uso por hispano-hablantes por Dairelys García Rivas en 2017-07

% INSTRUCCIONES:

% La idea es no editar mucho este documento, sino llenar la información
% en los diferentes ficheros en las carpetas;
% Configuraciones, Páginas_Princ, Capítulos, Anexos y posiblemente Figuras,
% así como el fichero Bibliografía.bib

% Esta plantilla es fácil de escalar a sus necesidades, 
% simplemente comente las declaraciones input que se explican debajo

% Configurar el documento:
\documentclass[letterpaper, 12pt, oneside]{Tesis}  % Usar el estilo Tesis", basado en el estilo ECS Thesis por Steve Gunn
\sffamily

% Adicionar más paquetes en Paquetes.tex:
% Add your own packages here, 
% these existing packages can be removed if necessary
% not that more packages are imported in Thesis.cls, '
%  but those should not be changed if you don't know what you are doing...


\usepackage[utf8]{inputenc} % for writing other that basic characters
\usepackage{graphicx}
\usepackage{caption}
\usepackage{subcaption}

\usepackage{pslatex}
\usepackage{pdfpages}

\usepackage{pdflscape}

% Include any extra LaTeX packages required
\usepackage[square, numbers, comma, sort&compress]{natbib}  % Use the "Natbib" style for the references in the Bibliography
\usepackage{verbatim}  % Needed for the "comment" environment to make LaTeX comments
%5%\usepackage{vector}  % Allows "\bvec{}" and "\buvec{}" for "blackboard" style bold vectors in maths

\usepackage{float} %Para que las imágenes se puedan poner donde desees

\usepackage{supertabular} %para dividir tablas que no caben en la página

% Usar si quiere:
%5%\graphicspath{Figuras/}  % Ruta del fichero de gráficos (configurado para que los gráficos estén en formato PDF)
%5%\hypersetup{urlcolor=blue, colorlinks=true}  % Colorea los hipervínculos en azul, pero puede distraer si hay muchos vínculos.
\hypersetup{urlcolor=black, colorlinks=true}

% Establecer su nombre, el título del trabajo y más en Administrativo.tex:
% Aquí es donde se definen el autor, la universidad, el título y más

% Información personal:
\newcommand{\myAuthorName}  {El autor de esta tesis}% Nombre del Autor
\newcommand{\myAuthorEmail} {micorreo@algo.com} % Correo electrónico del Autor
\newcommand{\myTutorOne} {El tutor principal}
\newcommand{\myTutorOneEmail} {correo@algo.com}
% Algunas tesis tienen más de un tutor
\newcommand{\myTutorTwo} {El segundo tutor}
\newcommand{\myTutorTwoEmail} {correo@algo.com}

\newcommand{\myTitle}       {EL TÍTULO DE LA TESIS} % Thesis title goes here
\newcommand{\mySubject}     {Tema de la tesis} % Subject goes here
\newcommand{\mySpKeywords}    {clave, palabras} % Spanish keywords goes here
\newcommand{\myKeywords}    {key, words} % Keywords goes here


% University information
\newcommand{\myUniversity}{Universidad Tecnol\'{o}gica de La Habana “Jos\'{e} Antonio Echeverr\'{i}a”} %The Iniversity name goes here
%\\ Facultad de Ingeniería Informática \\ Filial de Ciencias Técnicas de 10 de octubre
\newcommand{\myUniversityWeb}{} %University Web Site URL Here (include http://
\newcommand{\myDepartment}{Departamento del tesista} % The Department goes here 
\newcommand{\myDepartmentWeb}{} % Department Web Site URL Here (include http://)
\newcommand{\myGroup}{Grupo del tesista}
\newcommand{\myGroupWeb}{}
\newcommand{\myFaculty}{Facultad del tesista}
\newcommand{\myFacultyWeb}{}

%Degree, program or course: ex: Master of Science, Engineering Physics
\newcommand{\myDegree}{Ingenier\'{i}a en Inform\'{a}tica} % The degree, program or course-name goes here


% can be left untouched, both:
\newcommand{\myDate}{\today}
\newcommand{\myPartyalFulfillment}{Trabajo de diploma para optar por el t\'{i}tulo de }

% Nombres de los capitulos
\newcommand{\CapUno}{Título del capítulo Uno}
\newcommand{\CapDos}{Título del capítulo Dos}
\newcommand{\CapTres}{Título del capítulo Tres}
\newcommand{\CapCuatro}{Título del capítulo Cuatro}
\newcommand{\CapCinco}{Título del capítulo Cinco}

%% Para hacer el índice del Glosario de Términos
%\usepackage{makeidx}
%\makeindex

%% ----------------------------------------------------------------
\begin{document}
\frontmatter      % Comenzar estilo romano (i, ii, iii, iv...) para el numerado de páginas

%Márgenes
%\headheight 0in
\textheight = 25cm
\textwidth = 19cm % Ancho del texto
\topmargin = 1cm %Margen superior
\oddsidemargin = 2 cm %Margen izquierdo
%\parindent = 0mm % Sin sangría


% Aquí se importan las primeras páginas, puede encontrarlas en la carpeta Páginas_Princ
% Los ficheros en la subcarpeta Fijo no necesitan edición.
% Si no necesita ninguna de estas secciones, simplemente comente, o elimine la fila input

%% Todas las páginas antes de los capítulos ------------------------------
\begin{titlepage}

	\begin{center}
		\begin{figure}
			\includegraphics[width=17cm,height=1cm]{Figuras/separador.png}
		\end{figure}
		\sf {\large \textbf{\myUniversity}}
	
		{\large \textbf{\myFaculty}}
	
		{\large \textbf{\myGroup}}
	    
		\vspace{2em}
		
		\begin{figure}[h]
			\centering
			\includegraphics[width=3cm,height=4.5cm]{Figuras/Logo.png}
		\end{figure}
		
		\vspace{1em}

		{\large \textbf{\myTitle}}
		
		\vspace{1em}
		
		{\large \textbf{\myPartyalFulfillment}}
		{\large \textbf{\myDegree}}
		
	\end{center}
	
	\begin{table}[ht] 
		\centering 
		\begin{tabular}{ l  r  l }
			\hspace{6cm} & \sf \textbf{Autor:} & \sf \textbf{\myAuthorName} \\
			\hspace{3cm} &  & \sf \textbf{\myAuthorEmail} \\
			
			\hspace{3cm} & \sf \textbf{Tutores:} & \sf \textbf{\myTutorOne} \\
			\hspace{3cm} &  & \sf \textbf{\myTutorOneEmail} \\
			\hspace{3cm} &  & \sf \textbf{\myTutorTwo} \\
			\hspace{3cm} &  & \sf \textbf{\myTutorTwoEmail} \\
		\end{tabular}
	\end{table}
	
	\vspace{1em}	
	
	\begin{center}
		\sf \textbf{La Habana, Cuba}
		
		\sf \textbf{Mes Año}
		
		\begin{figure}
			\includegraphics[width=17cm,height=1cm]{Figuras/separador.png}
		\end{figure}
	\end{center}

\end{titlepage}

\thispagestyle{empty}

\clearpage

\vspace*{16 cm}

\begin{flushright}
\sf	

\textit{Frase interesante de algún personaje que haya inspirado al tesista.}

\textbf{Autor de la frase}
\end{flushright}

% Página de Agradecimientos, para agradecer a todos
\clearpage
\setstretch{1.3} % Resetear el espaciado a 1.3 para el cuerpo del texto (si cambia)
\addtocontents{toc}{\vspace{1em}}  % Incluir un espacio en los Contenidos, por estética
\acknowledgements{
	
	% Los agradecimientos y las personas van aquí, no olvidar incluir al supervisor del proyecto, tutores...
	
	\sf \paragraph{}Agradecimientos a todas las personas que el tesista considere que contribuyeron a la realización de la tesis.
	
}


\clearpage
\lhead{}  % Dejar vacío el encabezado superior izquierdo.
\setstretch{1.3}  % Devolver el espaciado a 1.3
\pagestyle{empty}  % El estilo de página necesita estar vacío para esta página

\dedicatory{
	\addtocontents{toc}{\vspace{1em}} % Añadir un espacio en los contenidos, por estética
	
	\vspace{140 mm}
	
	\begin{flushright}
		\sf {\Large \textbf{\emph{Dedicatoria a las personas\\
					que el tesista considere importantes.}}}
	\end{flushright}}



\clearpage 
\addtotoc{Síntesis}  % Adiciona la entrada "Síntesis" a los Contenidos

\thispagestyle{empty}
%\null\vfill
\begin{center}
	\setlength{\parskip}{0pt}
	\bigskip
	{\begin{flushleft}
			\sf \centering \textbf{\large{Resumen}}
		\end{flushleft} \par}
	\bigskip
\end{center}

\sf 
Resumen en español de la tesis.

\paragraph{Palabras clave:} clave, palabras.

% La Página Resumen (en inglés)
\clearpage 
\addtotoc{Abstract}  % Adiciona la entrada "Abstract" a los Contenidos
\addtocontents{toc}{\vspace{1em}}  % Incluir un espacio en los Contenidos, por estética

\abstract{
	
	\sf
	
\textit{	Resumen en idioma inglés de la tesis.}


\paragraph{Keywords:} key, words.

\clearpage % Comenzar una nueva página
\setstretch{1.3} % Reestablecer el espaciado a 1.3 para el cuerpo del texto (si cambia)
\pagestyle{fancy} % Los estilos de encabezado de página han estado "vacíos" todo este tiempo, 
                  % ahora se usan los encabezados "fancy" como se definió antes
\lhead{\emph{\'{I}ndice}}  % Cambiar el encabezado izquierdo superior a "Contenidos"
\setcounter{tocdepth}{3} % Para que tome hasta el tercer nivel de texto en los Contenidos
\renewcommand{\contentsname}{\'{I}ndice}
\tableofcontents  % Escribir la tabla de contenidos


\setstretch{1.3} % Reset the line-spacing to 1.3 for body text (if changed)
\pagestyle{fancy} % The page style headers have been "empty" all this time, 
                  % now use the "fancy" headers as defined before
\lhead{\emph{\'Indice de Figuras}}  % Cambiar el encabezado izquierdo superior a "Indice de Figuras"

\renewcommand{\listfigurename}{\'Indice de Figuras}
\listoffigures  % Write out the List of Figures


\clearpage  % Start a new page
\setstretch{1.3} % Reset the line-spacing to 1.3 for body text (if changed)
\pagestyle{fancy} % The page style headers have been "empty" all this time, 
                  % now use the "fancy" headers as defined before
\lhead{\emph{\'Indice de Tablas}}  % Cambiar el encabezado izquierdo superior a "Lista de Tablas"
\renewcommand{\listtablename}{\'Indice de Tablas}
\listoftables  % Write out the List of Tables


%% El Cuerpo -------------------------------------------------------
\setstretch{1.3}  % Devolver el espaciado de línea a 1.3
\mainmatter	  % Comenzar el numerado de página normal (1,2,3...)
\pagestyle{fancy}  % Devolver los encabezados de página al estilo "fancy"

\renewcommand{\chaptername}{Capítulo}
% Incluír los capítulos de la tesis, como ficheros separados
% Simplemente quite el símbolo de comentario a las líneas a medida
% que vaya escribiendo los capítulos

\addtocontents{toc}{\vspace{1em}} % Añadir un espacio en los contenidos, por estética

\lhead{\emph{Introducción}}  % Establecer el encabezado superior izquierdo a "Introducción"
% Para que no le ponga delante a la Introducción la palabra Capítulo
\chapter*{Introducción}
\addcontentsline{toc}{chapter}{Introducción}

\paragraph{}Lorem ipsum dolor sit amet, consectetur adipiscing elit, sed do eiusmod tempor incididunt ut labore et dolore magna aliqua. Ut enim ad minim veniam, quis nostrud exercitation ullamco laboris nisi ut aliquip ex ea commodo consequat. Duis aute irure dolor in reprehenderit in voluptate velit esse cillum dolore eu fugiat nulla pariatur. Excepteur sint occaecat cupidatat non proident, sunt in culpa qui officia deserunt mollit anim id est laborum \cite{noauthor_something_2007}.

Curabitur pretium tincidunt lacus. Nulla gravida orci a odio. Nullam varius, turpis et commodo pharetra, est eros bibendum elit, nec luctus magna felis sollicitudin mauris. Integer in mauris eu nibh euismod gravida. Duis ac tellus et risus vulputate vehicula. Donec lobortis risus a elit. Etiam tempor. Ut ullamcorper, ligula eu tempor congue, eros est euismod turpis, id tincidunt sapien risus a quam. Maecenas fermentum consequat mi. Donec fermentum. Pellentesque malesuada nulla a mi. Duis sapien sem, aliquet nec, commodo eget, consequat quis, neque. Aliquam faucibus, elit ut dictum aliquet, felis nisl adipiscing sapien, sed malesuada diam lacus eget erat. Cras mollis scelerisque nunc. Nullam arcu. Aliquam consequat. Curabitur augue lorem, dapibus quis, laoreet et, pretium ac, nisi. Aenean magna nisl, mollis quis, molestie eu, feugiat in, orci. In hac habitasse platea dictumst.

\subsubsection*{Situación Problemática:}
\paragraph{}Lorem ipsum dolor sit amet, consectetur adipiscing elit, sed do eiusmod tempor incididunt ut labore et dolore magna aliqua. Ut enim ad minim veniam, quis nostrud exercitation ullamco laboris nisi ut aliquip ex ea commodo consequat. Duis aute irure dolor in reprehenderit in voluptate velit esse cillum dolore eu fugiat nulla pariatur. Excepteur sint occaecat cupidatat non proident, sunt in culpa qui officia deserunt mollit anim id est laborum.

\subsubsection*{Problema a resolver:}
\paragraph{}¿Lorem ipsum dolor sit amet, consectetur adipiscing elit, sed do eiusmod tempor incididunt ut labore et dolore magna aliqua. Ut enim ad minim veniam, quis nostrud exercitation ullamco laboris nisi ut aliquip ex ea commodo consequat. Duis aute irure dolor in reprehenderit in voluptate velit esse cillum dolore eu fugiat nulla pariatur. Excepteur sint occaecat cupidatat non proident, sunt in culpa qui officia deserunt mollit anim id est laborum?
\paragraph{}El presente trabajo tiene como \textbf{objeto de estudio} en el \textbf{escenario del cliente/usuario}, ..... En el \textbf{escenario de la informática} se centra en .....
\paragraph{}Dentro del \textbf{campo de acción} en el \textbf{escenario del cliente/usuario} se encuentran ..... En el \textbf{escenario de la informática} se centra en .....
\paragraph{}Para darle solución a la problemática planteada se define como \textbf{objetivo general}, del presente trabajo de diploma, .....
\paragraph{}Para dar cumplimiento al objetivo general, será necesario dar cumplimiento a los \textbf{objetivos específicos} que se desglosan a continuación:
\begin{enumerate}
	\item Lorem ipsum dolor sit amet, consectetur adipiscing elit, sed do eiusmod tempor incididunt ut labore et dolore magna aliqua
\end{enumerate}

\paragraph{}Para dar cumplimiento a los diferentes objetivos específicos se establecieron las siguientes \textbf{tareas de investigación} por objetivo:

\textbf{Objetivo 1 ....}
\begin{enumerate}
	\item Hacer algo.

\subsubsection*{Valor Práctico}
\paragraph{}Con el desarrollo del sistema propuesto se espera, como aporte práctico: .....
\paragraph{}Para dar cumplimiento al objetivo propuesto el documento se estructura en cinco capítulos: \CapUno, \CapDos, \CapTres, \CapCuatro, \CapCinco. Cada cap\'{i}tulo está estructurado en introducción, desarrollo y conclusiones parciales.
\paragraph{Capítulo 1: \CapUno.} Explicación del capítulo 1.
\paragraph{Capítulo 2: \CapDos.} Explicación del capítulo 2. 
\paragraph{Capítulo 3: \CapTres.} Explicación del capítulo 3.
\paragraph{Capítulo 4: \CapCuatro.} Explicación del capítulo 4.
\paragraph{Capítulo 5: \CapCinco.} Explicación del capítulo 5.
\paragraph{}Por último, se presentan las conclusiones, recomendaciones, referencias bibliográficas, así como un glosario de términos que ayuda en el entendimiento de los conceptos específicos de la investigación. % Introducción

% Para todos los capítulos, use la funcionalidad nueva chap{} en vez de chapter{}
% Esto hará que el texto en la izquierda superior de la página sea el mismo que el del capítulo

% Los Nombres de los capítulos se definen en Configuraciones/Administrativo.tex
\chap{\CapUno}

\section{Introducción}
\sf \paragraph{}Lorem ipsum dolor sit amet, consectetur adipiscing elit, sed do eiusmod tempor incididunt ut labore et dolore magna aliqua. Ut enim ad minim veniam, quis nostrud exercitation ullamco laboris nisi ut aliquip ex ea commodo consequat. Duis aute irure dolor in reprehenderit in voluptate velit esse cillum dolore eu fugiat nulla pariatur. Excepteur sint occaecat cupidatat non proident, sunt in culpa qui officia deserunt mollit anim id est laborum.

\section{Sección Uno}
\paragraph{}Lorem ipsum dolor sit amet, consectetur adipiscing elit, sed do eiusmod tempor incididunt ut labore et dolore magna aliqua. Ut enim ad minim veniam, quis nostrud exercitation ullamco laboris nisi ut aliquip ex ea commodo consequat. Duis aute irure dolor in reprehenderit in voluptate velit esse cillum dolore eu fugiat nulla pariatur. Excepteur sint occaecat cupidatat non proident, sunt in culpa qui officia deserunt mollit anim id est laborum.

\subsection{Subsección uno}
\paragraph{}Lorem ipsum dolor sit amet, consectetur adipiscing elit, sed do eiusmod tempor incididunt ut labore et dolore magna aliqua. Ut enim ad minim veniam, quis nostrud exercitation ullamco laboris nisi ut aliquip ex ea commodo consequat. Duis aute irure dolor in reprehenderit in voluptate velit esse cillum dolore eu fugiat nulla pariatur. Excepteur sint occaecat cupidatat non proident, sunt in culpa qui officia deserunt mollit anim id est laborum. \ref{fig:afigure}

Lorem ipsum dolor sit amet, consectetur adipiscing elit, sed do eiusmod tempor incididunt ut labore et dolore magna aliqua. Ut enim ad minim veniam, quis nostrud exercitation ullamco laboris nisi ut aliquip ex ea commodo consequat. Duis aute irure dolor in reprehenderit in voluptate velit esse cillum dolore eu fugiat nulla pariatur. Excepteur sint occaecat cupidatat non proident, sunt in culpa qui officia deserunt mollit anim id est laborum. \ref{table:atable}

\paragraph{}Ejemplo de una tabla.

\begin{table}[ht] 
	\centering 
	\begin{tabular}{|p{3cm}|p{1.5cm}|p{1.5cm}|p{1.5cm}|p{1.5cm}|p{3cm}|p{3cm}| p{5cm}|}
		\hline
		\textbf{Criterio} & \textbf{Criterio} & \textbf{Criterio} & \textbf{Criterio} & \textbf{Criterio} & \textbf{Criterio} & \textbf{Criterio} \\ \hline
		Algo & Algo & Algo & Algo & Algo & Algo & Algo \\ \hline
		Algo & Algo & Algo & Algo & Algo & Algo & Algo \\
		\hline
	\end{tabular}
	\caption[Titular de la tabla]{Titular de la tabla}
	\label{table:atable}
\end{table}

\begin{figure}[H]
	\centering
	\includegraphics[width=6cm,height=6.5cm]{Figuras/Logo.png}
	\caption{\sf Título de la imagen.}
	\label{fig:afigure}
\end{figure}

\section{Conclusiones parciales}
\begin{itemize}
	\item Lorem ipsum dolor sit amet, consectetur adipiscing elit, sed do eiusmod tempor incididunt ut labore et dolore magna aliqua. Ut enim ad minim veniam, quis nostrud exercitation ullamco laboris nisi ut aliquip ex ea commodo consequat. Duis aute irure dolor in reprehenderit in voluptate velit esse cillum dolore eu fugiat nulla pariatur. Excepteur sint occaecat cupidatat non proident, sunt in culpa qui officia deserunt mollit anim id est laborum.
\end{itemize} % 

% % Para todos los capítulos, use la funcionalidad nueva chap{} en vez de chapter{}
% Esto hará que el texto en la izquierda superior de la página sea el mismo que el del capítulo

% Los Nombres de los capítulos se definen en Configuraciones/Administrativo.tex
\chap{\CapDos} 

\section{Introducción}
\paragraph{}En este capítulo se hace referencia al estudio para realizar la descripción detallada sobre el modelo del negocio, del cual se obtendrá posteriormente la captura de requisitos y la modelación. Además, se definen los actores y trabajadores que intervienen en el negocio, así como las reglas y se exponen sus respectivas descripciones de acuerdo con los roles que desempeñan. Mediante el uso del lenguaje UML se representa el diagrama de casos de uso del negocio, los diagramas de actividades de los casos de uso definidos y el modelo de objetos que muestra las relaciones entre trabajadores y entidades. Se enumeran los requisitos funcionales y se muestra una breve descripción de los no funcionales que se satisfacen en la construcción de la solución propuesta. Se realiza una descripción de los actores y los casos de uso del sistema. Además, se identifican los paquetes a implementar y la relación existente entre ellos, con el objetivo de administrar la complejidad.

\section{Modelo del negocio actual}
\paragraph{}El negocio se centra en la oficina de la caja, local donde se registran sus operaciones. El local se encuentra ubicado físicamente en el segundo piso del edificio central del CIM. Este contiene la seguridad necesaria para la custodia del efectivo la cual se garantiza a través de cámaras y personal entrenado. El cajero tiene acceso a esta oficina, dentro se encuentran dos cajas fuertes para la custodia de todo el efectivo. La combinación de ambas cajas es de conocimiento únicamente del cajero, existiendo una copia lacrada en un sobre que se ubica en la direcci\'{o}n del CIM.
\paragraph{}El negocio consta de dos tipos de acciones fundamentales, que generan operaciones de soporte y completamiento: la entrega y la recepción de efectivo. 
\paragraph{}La recepción de efectivo en la caja consiste en el acto de recibir billetes y monedas por parte del cajero, por los siguientes motivos: venta de tickets en el comedor, recogida de pasaje de ómnibus, el cobro de efectivo del servicio de lavandería, la venta de medicamentos en el consultorio, el cobro por concepto de pérdida del solapín, entre otros. Estas operaciones son realizadas fundamentalmente por los jefes de las áreas de la UEB que atiende los servicios en el CIM y que se llama SERVICIM. Por cada entrega de efectivo, el cajero debe confeccionar un modelo denominado Recibo de Efectivo, después de verificar que los justificantes y el efectivo que traen los jefes de área coincidan. 
\paragraph{}Por la suma de varios recibos de efectivo, el cajero realiza en el banco un depósito. La entrega del efectivo por parte del cajero ocurre cuando los trabajadores de las áreas del CIM: necesitan viajar a provincia por cuestiones de trabajo, o tienen que realizar pagos menores (ponches, confección de llaves, reparaciones menores, compra de sellos para realizar legalizaciones, entre otros) que no excedan la cantidad establecida en la caja. Esta actividad se centraliza por los jefes de las áreas y por cada operación el cajero emite un modelo de vale para pagos menores por los justificantes entregados. Para el caso de los viajes, se confecciona un modelo de Anticipos y Liquidación de Gastos de Viaje.
\paragraph{}La caja tiene un fondo límite para operar y en ningún caso se puede utilizar el efectivo ingresado para realizar pagos. Es por esto que por la suma de los Vales para Pagos Menores y los modelos de dietas emitidas, el cajero debe reponer el efectivo entregado mediante la confección de un modelo de Reembolso del Fondo para Pagos Menores.
\paragraph{}El efectivo en la caja es controlado a través de arqueos sorpresivos por parte del jefe del Departamento de Finanza (el cual realiza el del último día del mes para fijar los saldos en la contabilidad). El efectivo es contado y al saldo obtenido se le suman los vales de pagos menores y las dietas que no han podido ser reembolsadas. Se verifica que la suma de ambos importes no exceda del límite monetario aprobado para estar en poder del cajero.
\paragraph{}Los modelos de anticipos de dietas se controlan por parte del especialista de contabilidad en un modelo denominado Control de los Anticipos a Justificar, por el cual se lleva la trazabilidad de las dietas, sus liquidaciones para su contabilización correspondiente.
\paragraph{}La caja además realiza otras operaciones, como la recepción de cheques por parte del especialista de contabilidad, de pagos realizados por otras empresas, los cuales se depositan en la cuenta bancaria. Igualmente se reciben sellos contra el vale entregado para estos fines.
\paragraph{}Todas estas operaciones deben ser anotadas por el cajero en el modelo Control Diario de Operaciones de Caja.

\section{Reglas del negocio a considerar}
\begin{itemize}
	\item El efectivo por vale extraído de la caja no puede exceder los 500.00 CUP, ni los 500.00 CUC.
	\item Todos los modelos deben ser firmados por personal autorizado para ello. El cual debe ser trabajador del CIM.
	\item La suma captada de los montos autorizados no puede ser superior al fondo aprobado.
	\item La operaciones de Anticipos para Gastos de Viaje no pueden exceder los montos autorizados por tipo de operaci\'{o}n.
	\item El efectivo existente en la caja no puede exceder los 10,000.00 CUP, ni los 1,000.00 CUC (fondo autorizado para cada tipo de moneda).
	\item Todos los modelos tienen que tener su desglose monetario en pesos y centavos por tipo de moneda.
	\item Los documentos que amparen entrega y/o recepci\'{o}n de efectivo tienen una vigencia m\'{a}xima de 5 años.
	\item Los documentos no pueden aparecer firmados en escaques diferentes por la misma persona y las firmas deben ser de personas autorizadas.
	\item El anticipo para compra de sellos debe realizarse en un modelo de Vale para Pagos Menores Provisional, el que será cancelado una vez que el usuario 	autorizado entregue los sellos, para emitir un modelo definitivo con el 	desglose correspondiente.
	\item Para realizar el arqueo de caja deben sumarse documentos y contarse los 	valores por denominación, la suma de ambos se compara con el fondo autorizado por tipo de moneda.
	\item El combustible por tarjeta magnética tendrá el tratamiento según la norma contable, se emitirán cheques por el valor de la asignación de combustible y	la compra de tarjetas nuevas en la moneda que se designe. La carga se hace por listado de asignación de la empresa. El gasto se refleja a través de Vales	para Pagos Menores. Existe software en el Departamento de Transporte para	determinar vencimientos y otros parámetros.
	\item Los modelos utilizados para amparar salidas y entradas de efectivo de la caja deben tener el formato establecido en la norma vigente \cite{noauthor_resolucion_2007}.
	\item Los modelos deben tener numeración y esta es independiente por tipo de 	moneda: CUC y CUP, ambos deben cumplir los requisitos de la normativa vigente.
	\item El efectivo que se reciba en la caja debe ser depositado en banco, los pagos menores se realizarán del fondo disponible.
\end{itemize}

\section{Actores del negocio}
\begin{table}[ht]
	\begin{supertabular}{ |p{6cm}|p{10cm}| }
		\hline
		\textbf{Nombre del actor}
		& \textbf{Descripción} \\ \hline
		
		Jefe de Departamento Finanzas
		& Es la persona encargada de realizar las supervisiones al efectivo mediante arqueos sorpresivos. Revisa las operaciones realizadas. \\ \hline
		
		Cliente
		& Representa a los jefes de las áreas o personal autorizado que se presente en la caja para realizar una operación que involucre efectivo, sellos, cheques o monedas. \\
		
		\hline
	\end{supertabular}
	\caption[Actores del Negocio]{Actores del Negocio}
	\label{table:act_neg}
\end{table}

\section{Diagrama de casos de uso del negocio}
\paragraph{}El negocio contará con dos casos de uso, los cuales se grafican en el diagrama que se muestra a continuación en las Figuras \ref{fig:CUN_1} y \ref{fig:CUN_2}.

\begin{figure}[H] %Con el paquete float, se pone la opción [H] para que te salga la imagen donde la quieres
	\centering
	\includegraphics[width=9cm,height=5cm]{Figuras/CUN_1.jpg}
	\caption{Caso de uso del actor Jefe de Departamento de Finanzas}
	\label{fig:CUN_1}
\end{figure}

\begin{figure}[H] %Con el paquete float, se pone la opción [H] para que te salga la imagen donde la quieres
	\centering
	\includegraphics[width=9cm,height=3cm]{Figuras/CUN_2.jpg}
	\caption{Caso de uso del actor Cliente}
	\label{fig:CUN_2}
\end{figure}

\section{Trabajadores del negocio}
\begin{table}[ht]
	\begin{supertabular}{ |l|p{12cm}| }
		\hline
		\textbf{Nombre del trabajador}
		& \textbf{Descripción} \\ \hline
		
		Cajero
		& Es la persona encargada de realizar todas las operaciones solicitadas por los clientes y el jefe de departamento. Tiene control sobre el efectivo. \\ \hline
		
		Especialista de contabilidad
		& Es la persona encargada de recibir los registros de cierre diario de las operaciones de caja. \\
		
		\hline
	\end{supertabular}
	\caption[Trabajadores del Negocio]{Trabajadores del Negocio}
	\label{table:trb_neg}
\end{table}

\section{Casos de uso del negocio}
\paragraph{}La secuencia de acciones de un caso de uso del negocio se describe mediante un flujo de trabajo y se puede representar mediante una descripción literal del caso de uso o mediante el diagrama de actividades \cite{miranda_aplicaciones_2018}.
\paragraph{}A continuación, se describen los casos de uso del negocio a través de los diagramas de actividades.

\subsection{Caso de Uso 1: Solicitar cierre Control Diario de Operaciones de Caja.}
\begin{figure}[H] %Con el paquete float, se pone la opción [H] para que te salga la imagen donde la quieres
	\centering
	\includegraphics[width=19cm,height=18cm]{Figuras/CU_1.jpg}
	\caption{Diagrama de Actividades CUN Solicitar cierre diario de operaciones de caja.}
	\label{fig:CU_1}
\end{figure}

\subsection{Caso de Uso 2: Realizar arqueo de caja}
\begin{center}
	\begin{figure}[H] %Con el paquete float, se pone la opción [H] para que te salga la imagen donde la quieres
		\centering
		\includegraphics[width=18cm,height=17cm]{Figuras/DACUNRArqueo.jpg}
		\caption{\sf Diagrama de Actividades CUN Realizar arqueo de caja.}
		\label{fig:DACUNRArqueo}
	\end{figure}
\end{center}


\subsection{Caso de Uso 3: Solicitar operación}
\begin{figure}[H] %Con el paquete float, se pone la opción [H] para que te salga la imagen donde la quieres
	\centering
	\includegraphics[width=17cm,height=20cm]{Figuras/DACUNSolOp.png}
	\caption{\sf Diagrama de Actividades CUN Solicitar operación.}
	\label{fig:DACUNSolOp}
\end{figure}

\section{Modelo de objetos}
\begin{figure}[H] %Con el paquete float, se pone la opción [H] para que te salga la imagen donde la quieres
	\centering
	\includegraphics[width=14cm,height=10cm]{Figuras/Mod_Obj.jpg}
	\caption{Diagrama de objetos del negocio.}
	\label{fig:Mod_Obj}
\end{figure}

\section{Conclusiones parciales}
\begin{itemize}
	\item Se obtuvieron los	casos de uso del negocio que permiten definir las entidades.
	\item Se definieron como trabajadores del negocio al cajero y al jefe de departamento de contabilidad, lo que permite establecerlos como actores del sistema.
	\item El trabajador “Cajero” es el que interactúa con todas las entidades.
\end{itemize} % 

% % Para todos los capítulos, use la funcionalidad nueva chap{} en vez de chapter{}
% Esto hará que el texto en la izquierda superior de la página sea el mismo que el del capítulo

% Los Nombres de los capítulos se definen en Configuraciones/Administrativo.tex
\chap{\CapTres}
\paragraph{}En el presente capítulo se explica de forma general el funcionamiento del sistema. Se enumeran los requisitos funcionales y se muestra una breve descripción de los no funcionales que se satisfacen en la construcción de la solución propuesta. Se realiza una descripción de los actores y los casos de uso del sistema. Además, se identifican los paquetes a implementar y la relación existente entre ellos, con el objetivo de administrar la complejidad.

\section{Actores del sistema a informatizar}
\paragraph{}En la tabla \ref{table:act_sist} se muestra la definición de los actores del sistema a informatizar.

\begin{table}%[H]
	\begin{supertabular}{|p{5cm}|p{12cm}|}
		\hline
		\textbf{Nombre del actor}
		& \textbf{Descripción} \\ \hline
		
		Administrador
		& Actor que tiene acceso a la creación de usuarios y verifica la traza del sistema así como la gestión de la salva y restaura de la base de datos en el servidor y la restaura en el sistema. \\ \hline
		
		Usuario
		& Actor genérico que tiene acceso a opciones del sistema relacionadas con, autenticar usuario, cambiar contraseña y visualizar ayuda y algunos de los reportes del sistema. \\ \hline
		
		Cajero
		& Actor encargado de crear los documentos en el sistema, Vales para Pagos Menores, Recibos de Efectivo, Anticipos y Liquidación de Gastos de Viaje, Depósito de Efectivo y Sellos, Movimiento Control Diario de Operaciones de Caja. \\ \hline
		
		Jefe de Departamento
		& Actor encargado de la gestión de nomencladores, tiene acceso al Arqueo de Caja para realizar revisiones sorpresivas, así como fijar en la contabilidad los saldos de cierre de mes y revisa diversos documentos entre los que se encuentran: Reembolso del Fondo para Dietas y Pagos Menores y Depósitos de Efectivo y Cheques, tiene acceso a la salva del sistema, realiza la exportación de los comprobantes de operaciones. \\ \hline
		
		Especialista de Contabilidad
		& Entrega los Cheques recibidos para su posterior depósito. Revisa para contabilizar los Depósitos de Efectivo y Cheques y los Reembolso del Fondo para Dietas y Pagos Menor, lleva el control de Anticipos a Justificar, la impresión de comprobantes de operaciones. \\
		
		\hline
	\end{supertabular}
	\caption[Actores del sistema a informatizar]{Actores del sistema a informatizar}
	\label{table:act_sist}
\end{table}

\section{Definición de los requisitos funcionales}
\paragraph{}Los requerimientos funcionales se definen como condiciones que deben cumplir un sistema para satisfacer un contrato, una norma, o sea, describen la funcionalidad o los servicios que se espera el sistema proveerán, sus entradas y sus salidas \cite{miranda_aplicaciones_2018}. A continuación, se exponen los requisitos funcionales de la aplicación empaquetados.
\begin{enumerate}
	\item Seguridad
	\begin{enumerate}
		\item Iniciar sesión de usuario
		\item Definir permisos por rol
		\item Asociar rol a usuario
		\item Cambiar contraseña
		\item Cambiar de usuario
		\item Visualizar trazas de usuarios
		\item Restaurar salva de la base de datos
		\item Realizar salva de la base de datos
	\end{enumerate}
	\item Configuración
	\begin{enumerate}
		\item Obtener lista de empleados
		\item Gestionar nomencladores
		\begin{enumerate}
			\item Provincias
			\item Municipios
			\item Cuentas contables\footnote{Se anotarán solamente, del nomenclador de cuentas contables emitido por la OSDE, aquellas que tienen que ver 	con la actividad}
			\item Conceptos de documentos a procesar\footnote{Estos conceptos los define la entidad, de acuerdo a su nivel de actividad y sirven para 	homogenizar el tratamiento de los documentos emitidos}
		\end{enumerate}
		\item Modificar fondo autorizado para pagos menores
		\item Modificar montos autorizados para gastos de viajes
		\item Insertar saldos iniciales desde documento de arqueo de caja
		\begin{enumerate}
			\item Introducir desglose de efectivo por denominación y moneda
		\end{enumerate}
	\end{enumerate}
	\item Operaciones
	\begin{enumerate}
		\item Efectuar pagos menores
		\item Efectuar Recibo de Efectivo
		\item Efectuar anticipo y liquidación de gastos de viaje
		\begin{enumerate}
			\item Calcular efectivo por día de viaje
		\end{enumerate}
		\item Efectuar reembolso del fondo para dietas y pagos menores
		\item Efectuar arqueo de caja
		\begin{enumerate}
			\item Calcular efectivo existente por movimientos realizados
		\end{enumerate}
		\item Efectuar depósito de efectivo y cheques
		\item Registrar cheque
		\item Gestionar recibo de sellos del timbre
		\item Obtener ficheros xml para exportar al sistema contable
	\end{enumerate}
	\item Reportes
	\begin{enumerate}
		\item Control de anticipos a justificar
		\item Reembolso del Fondo para Pagos Menores
		\item Movimiento Control Diario de Operaciones de Caja
		\item Movimientos mensuales
		\item Obtener ayuda sobre funcionalidades del sistema
		\item Obtener comprobantes de operaciones exportados al sistema contable
	\end{enumerate}
\end{enumerate}

\section{Paquetes y sus relaciones}
\paragraph{}Como se aprecia en la Figura \ref{fig:Paquetes}, se han agrupado las funcionalidades de la aplicación en 4 paquetes para facilitar el desarrollo de la misma.

\begin{figure}[H] %Con el paquete float, se pone la opción [H] para que te salga la imagen donde la quieres
	\centering
	\includegraphics[width=12cm,height=6cm]{Figuras/paquetes.jpg}
	\caption{Diagrama de paquetes de la aplicación}
	\label{fig:Paquetes}
\end{figure}

\paragraph{Paquete Seguridad:} En este paquete están los casos de uso asociados a la seguridad del sistema, contiene las funcionalidades relacionadas con los usuarios, contraseña, roles y visualización de las trazas del sistema.
\paragraph{Paquete Configuración:} Se encuentran aquí los casos de uso relacionados con la gestión de nomencladores del sistema como son: Usuarios, Empleados, Provincias, Municipios, Cuentas Contables, así como los conceptos de las operaciones a realizar la configuración general: Fondo autorizado para pagos menores y gastos de viaje y la introducción de los saldos iniciales para comenzar a explotar la aplicación.
\paragraph{Paquete Operaciones:} Se agrupan las funcionalidades principales del sistema, se encuentran los casos de uso: Arqueo de Caja, Reembolso del Fondo para Pagos Menores y Dietas, Vales para Pagos Menores, Control de Anticipos a Justificar, Recibo de Efectivo, Anticipo y Liquidación de Gastos de Viajes, Depósito de Efectivo y Cheques, Movimiento Control Diario de Operaciones de Caja y Entrega de Sellos, Emitir e imprimir comprobante de operaciones.
\paragraph{Paquete Reportes:} Se agrupan aquí los casos de uso de los reportes que se muestran en el sistema, así como con la visualización de la ayuda.

\section{Diagrama de Casos de Uso del sistema a informatizar}
\paragraph{}“Cada forma en que los actores usan el sistema se representa con un caso de uso. Los casos de uso son “fragmentos” de funcionalidad que el sistema ofrece para aportar un resultado de valor para sus actores. De manera más precisa, un caso de uso especifica una secuencia de acciones que el sistema puede llevar a cabo al interactuar con sus actores, incluyendo alternativas dentro de la secuencia. Por tanto, un caso de uso es una especificación. Especifica el comportamiento de “cosas” dinámicas, en este caso, de instancias de los casos de uso” \cite{goran_v._rational_2007}.
\paragraph{}En el Diagrama de Casos de Uso del Paquete Seguridad (ver Anexos, figura \ref{fig:DCU_Seg}). El rol administrador tendrá acceso a cinco de las opciones del sistema en este paquete, si está previamente autenticado para poder realizarlas.
\paragraph{}En el Diagrama Casos de Uso del Paquete Configuración (ver Anexos, figura \ref{fig:DCU_Conf}) el rol Jefe de Departamento de Finanzas tiene a su cargo la realización de seis funcionalidades en este paquete. Para ello tiene que estar autenticado en el sistema. 
\paragraph{}En el Diagrama de Casos de Uso del Paquete Operaciones (ver Anexo, figura \ref{fig:DCU_Oper}), en este paquete se muestran las interacciones de tres usuarios con un total de doce casos de usos repartidos entre ellos. Tres para Especialista en Contabilidad, dos para el Jefe de Departamento y siete para el Cajero. Para que un actor del sistema pueda acceder a sus funcionalidades, debe estar autenticado en la aplicación.
\paragraph{}En el Diagrama de Casos de Uso del Paquete Reportes (ver Anexos, figura \ref{fig:DCU_Rep}). En este paquete el rol de Usuario cuenta con las funcionalidades de Visualizar Reportes y Visualizar Ayuda, Cambiar Contraseña y Cambiar de Usuario, para ello tiene que estar autenticado en el sistema.

\section{Definición de los requisitos no funcionales}
\paragraph{}Los requisitos no funcionales son propiedades o cualidades que el producto de \textit{software} debe poseer \cite{miranda_aplicaciones_2018}. A continuación, se describen los requisitos no funcionales a tener en cuenta en el desarrollo del sistema propuesto.

\subsection{Apariencia o interfaz externa}
\begin{enumerate}
	\item El sistema tendrá una apariencia parecida a la del sistema contable en función de facilitar su uso, el sistema de ventanas está diseñado con igual propósito
	\item Se requiere el uso de un lenguaje familiar para el usuario 
	\item La interfaz principal debe contener la imagen del CIM, así como los colores que el Centro tiene definido en su manual de identidad corporativa
	\item Los modelos impresos tendrán el logotipo del CIM en su esquina superior izquierda
	\item Los modelos deben tener la estructura de la norma específica  del Ministerio de Finanzas y Precios\cite{noauthor_resolucion_2007}
\end{enumerate}
\subsection{Usabilidad}
\begin{enumerate}
	\item El sistema debe brindar facilidad en su uso con menús desplegables claros y espaciados adecuadamente, permitiendo una lectura fácil al usuario
	\item Se emplea un diseño similar en todas las interfaces, por lo que en todos los formularios se debe utilizar igual tipo de letra y deben contener un título en su parte superior izquierda que indique el nombre de la funcionalidad del sistema
	\item Los usuarios recibirán un adiestramiento que les permita entrenarse con la aplicación
	\item El sistema contará con ayuda, lo que permitirá al usuario independencia después de un entrenamiento básico
\end{enumerate}
\subsection{Seguridad}
\begin{enumerate}
	\item Sólo el personal claramente identificado y autorizado tendrá acceso a sistemas en funcionamiento y a los datos \cite{jacobson_casos_2013}
	\item Para la autenticación, la contraseña tendrá una longitud entre cuatro y seis caracteres y será cifrada, la misma tiene que ser una contraseña segura compuesta de letras y números
	\item Se realizará chequeo de seguridad sobre las acciones realizadas mediante un registro de trazas en el servidor
	\item El sistema brindará la posibilidad de salvar y restaurar la base de datos de la aplicación
	\item El sistema exportará las operaciones en ficheros de tipo XML, compatibles con el sistema contable Siscont V disponible en CIM, en caso que la importación falle, pueden teclearse directamente
	\item Los usuarios de las tecnologías de la información están en la obligación de informar de inmediato cualquier incidente de seguridad, debilidad o amenaza a sistemas o servicios y las direcciones correspondientes exigirán su cumplimiento \cite{jacobson_casos_2013}
\end{enumerate}

\subsection{Confiabilidad}
\begin{enumerate}
	\item La información manejada por el sistema estará protegida de acceso no autorizado mediante la autenticación de los usuarios con encriptación de contraseña, para evitar que ésta pueda ser descifrada por intrusos 
	\item Cada usuario solamente podrá acceder a las funcionalidades permitidas según su rol
	\item El Administrador del sistema será el usuario encargado de otorgar permisos, manejar roles, consultar la bitácora y restaurar la base de datos
	\item La base de datos del sistema estará ubicada en un servidor corporativo
	\item La base de datos tendrá un plan de salvas programadas automáticamente en el servidor, así como la posibilidad de salvar por parte de los usuarios del sistema
	\item El Jefe de Departamento de Finanzas será la persona encargada de la salva	de las bases de datos del sistema y la misma estará ubicada en la carpeta	MSQLBACKUP del disco de la partición D del disco duro. El administrador	del sistema será el único usuario con posibilidades de restaurar la base de	datos
\end{enumerate}
\subsection{\textit{Software}}
\begin{enumerate}
	\item Las computadoras de los usuarios que empleen la aplicación deben tener instalado el sistema operativo \textit{Windows} XP como mínimo, o Linux
	\item Igualmente en las estaciones deberá estar instalada la Máquina Virtual de Java 1.7.0
	\item El servidor de base de datos debe tener instalado el sistema operativo \textit{Windows Server} 2000 o superior
	\item Para que la aplicación funcione se utilizará un SGBD en un servidor	corporativo del CIM que se encuentra en el Departamento de Informática y	contiene la base de datos distribuida a la que se conecta el sistema
	\item El servidor debe ser multiusuarios, multiplataforma, tener soporte para almacenar como mínimo las computadoras de los usuarios que empleen grandes bases de datos y consumir pocos recursos
\end{enumerate}
\subsection{\textit{Hardware}}
\paragraph{}Como mínimo las computadoras de los usuarios que empleen la aplicación deben contar con 1.8 GHz de microprocesador, 512 Mb de RAM, 1 Gigabyte disponible en disco duro.
\begin{enumerate}
	\item La aplicación debe contar con 1.8 GHz de microprocesador, 512 Mb de RAM, 1 Gigabyte disponible en disco duro
	\item Como mínimo se debe disponer de un servidor de base de datos que cuente con 3.0 GHz de microprocesador, 1 Gigabyte de RAM y 10 Gigabytes disponibles en disco duro
	\item Al menos una impresora para imprimir los reportes
\end{enumerate}

\section{Descripción de los casos de uso}
\paragraph{}Para realizar la descripción se toma como base el listado de los casos de uso del sistema, al clasificarse los mismos de acuerdo a su impacto en la operatoria del sistema en crítico a aquellos casos de uso que se identifican como imprescindibles, secundario para aquellos que tienen un impacto medio en la aplicación y auxiliar para aquellos que sirven de base a otros casos de uso o que su impacto en el sistema sea menor \cite{jacobson_casos_2013}.

\begin{table}[H]
	\sf
	\begin{supertabular}{|p{4.5cm}|p{12cm}|}
		\hline
		\textbf{Nombre del caso de uso}
		& \textbf{Gestionar usuarios y roles} \\ \hline
		
		Actor
		& Administrador \\ \hline
		
		Resumen
		& El caso de uso inicia cuando el Administrador del sistema desea realizar la gestión de usuarios y sus roles correspondientes \\ \hline
		
		Precondiciones
		& El usuario debe estar creado en la base de datos del sistema el usuario debe estar en la tabla de empleado para poder acceder al sistema. \\ \hline
		
		Poscondiciones
		& El sistema debe mostrar los usuarios gestionados, así como su información y el rol que le pertenece en el sistema. \\ \hline
		
		Referencia
		& Paquete Seguridad. RF 1. \\ \hline

		Prioridad
		& Crítico \\		
		\hline
	\end{supertabular}
	\caption[Descripción del caso de uso Gestionar usuarios y roles]{Descripción del caso de uso Gestionar usuarios y roles}
	\label{table:GestUsers}
\end{table}

\begin{table}[H]
	\sf
	\begin{supertabular}{|p{4.5cm}|p{12cm}|}
		\hline
		\textbf{Nombre del caso de uso}
		& \textbf{Cambiar contraseña} \\ \hline
		
		Actor
		& Usuario \\ \hline
		
		Resumen
		& El caso de uso se inicia cuando el actor especificado desea cambiar su contraseña en el sistema. El caso de uso finaliza al quedar actualizada la contraseña \\ \hline
		
		Precondiciones
		& El usuario debe estar creado en la base de datos del sistema \\ \hline
		
		Poscondiciones
		& La nueva contraseña debe quedar guardada en la base de datos
		El usuario queda autenticado en el sistema \\ \hline
		
		Referencia
		& Paquete Seguridad. RF 6. \\ \hline
		
		Prioridad
		& Auxiliar \\		
		\hline
	\end{supertabular}
	\caption[Descripción del caso de uso Cambiar contraseña]{Descripción del caso de uso Cambiar contraseña}
	\label{table:CU_CambPass}
\end{table}

\begin{table}[H]
	\sf
	\begin{supertabular}{|p{4.5cm}|p{12cm}|}
		\hline
		\textbf{Nombre del caso de uso}
		& \textbf{Cambiar usuario} \\ \hline
		
		Actor
		& Usuario \\ \hline
		
		Resumen
		& El caso de uso se inicia cuando el actor especificado desea cambiar de usuario porque otro ocupará la estación de trabajo. El caso de uso finaliza cuando el usuario actual sale del sistema y aparece la ventana de cambio de usuario \\ \hline
		
		Precondiciones
		& El usuario debe estar creado en la base de datos del sistema \\ \hline
		
		Poscondiciones
		& El usuario actual cuando finalice la sesión, el sistema debe salir de esta \\ \hline
		
		Referencia
		& Paquete Seguridad. RF 7 \\ \hline
		
		Prioridad
		& Auxiliar \\		
		\hline
	\end{supertabular}
	\caption[Descripción del caso de uso Cambiar de usuario]{Descripción del caso de uso Cambiar de usuario}
	\label{table:CU_CambUser}
\end{table}

\begin{table}[H]
	\sf
	\begin{supertabular}{|p{4.5cm}|p{12cm}|}
		\hline
		\textbf{Nombre del caso de uso}
		& \textbf{Gestionar provincia} \\ \hline
		
		Actor
		& Jefe de Departamento de Contabilidad \\ \hline
		
		Resumen
		& El caso de uso se inicia cuando el actor especificado desea gestionar las provincias. El caso de uso permite insertar, modificar y eliminar provincias. El caso de uso termina con un listado de las provincias gestionadas \\ \hline
		
		Precondiciones
		& El usuario debe estar creado en la base de datos y estar autenticado en el sistema \\ \hline
		
		Poscondiciones
		& La información asociada a las provincias debe quedar actualizada en la base de datos \\ \hline
		
		Referencia
		& Paquete Configuración. RF 2.2 \\ \hline
		
		Prioridad
		& Secundario \\		
		\hline
	\end{supertabular}
	\caption[Descripción del caso de uso Gestionar provincia]{Descripción del caso de uso Gestionar provincia}
	\label{table:CU_GestProv}
\end{table}

\begin{table}[H]
	\sf
	\begin{supertabular}{|p{4.5cm}|p{12cm}|}
		\hline
		\textbf{Nombre del caso de uso}
		& \textbf{Gestionar municipio} \\ \hline
		
		Actor
		& Jefe de Departamento de Contabilidad \\ \hline
		
		Resumen
		& El caso de uso se inicia cuando el actor especificado desea gestionar los municipios. El caso de uso permite insertar, modificar y eliminar municipios. El caso de uso termina con un listado de los municipios gestionados \\ \hline
		
		Precondiciones
		& El usuario debe estar creado en la base de datos y estar autenticado en el sistema. Para crear un municipio deben estar creada la provincia a la que pertenece \\ \hline
		
		Poscondiciones
		& La información asociada a las provincias debe quedar actualizada en la base de datos \\ \hline
		
		Referencia
		& Paquete Configuración. RF 2.2 \\ \hline
		
		Prioridad
		& Secundario \\		
		\hline
	\end{supertabular}
	\caption[Descripción del caso de uso Gestionar municipio]{Descripción del caso de uso Gestionar municipio}
	\label{table:CU_GestMuni}
\end{table}

\begin{table}[H]
	\sf
	\begin{supertabular}{|p{4.5cm}|p{12cm}|}
		\hline
		\textbf{Nombre del caso de uso}
		& \textbf{Gestionar cuenta contable} \\ \hline
		
		Actor
		& Jefe de Departamento de Contabilidad \\ \hline
		
		Resumen
		& El caso de uso se inicia cuando el actor especificado desea gestionar las cuentas contables. El caso de uso permite insertar, modificar y eliminar cuentas contables. El caso de uso termina con un listado de las cuentas contables gestionadas \\ \hline
		
		Precondiciones
		& El usuario debe estar creado en la base de datos y estar autenticado en el sistema \\ \hline
		
		Poscondiciones
		& La información asociada a las cuentas contables debe quedar actualizada en la base de datos \\ \hline
		
		Referencia
		& Paquete Configuración. RF 2.2 \\ \hline
		
		Prioridad
		& Secundario \\		
		\hline
	\end{supertabular}
	\caption[Descripción del caso de uso Gestionar cuenta contable]{Descripción del caso de uso Gestionar cuenta contable}
	\label{table:CU_GestCCont}
\end{table}

\begin{table}[H]
	\sf
	\begin{supertabular}{|p{4.5cm}|p{12cm}|}
		\hline
		\textbf{Nombre del caso de uso}
		& \textbf{Obtener lista de empleados} \\ \hline
		
		Actor
		& Administrador \\ \hline
		
		Resumen
		& El caso de uso se inicia cuando el actor especificado desea gestionar los empleados del sistema. Para lo cual deberá importarlos de la base de datos  \\ \hline
		
		Precondiciones
		& El actor deberá tener acceso a la base de datos con la información requerida de los empleados \\ \hline
		
		Poscondiciones
		& La información asociada a los empleados debe quedar almacenada en la base de datos \\ \hline
		
		Referencia
		& Paquete Configuración. RF 2.1 \\ \hline
		
		Prioridad
		& Secundario \\		
		\hline
	\end{supertabular}
	\caption[Descripción del caso de uso Obtener lista de empleados]{Descripción del caso de uso Obtener lista de empleados}
	\label{table:CU_GetEList}
\end{table}

\begin{table}[H]
	\sf
	\begin{supertabular}{|p{4.5cm}|p{12cm}|}
		\hline
		\textbf{Nombre del caso de uso}
		& \textbf{Modificar fondo autorizado para pagos menores} \\ \hline
		
		Actor
		& Jefe de Departamento \\ \hline
		
		Resumen
		& El caso de uso se inicia cuando el actor especificado desea fijar el fondo máximo aprobado según documento firmado por el director de la Empresa. El caso de uso termina cuando el fondo ha sido fijado satisfactoriamente \\ \hline
		
		Precondiciones
		& El usuario tiene que haberse autenticado en el sistema y contar con los privilegios necesarios y contar con la información necesaria para realizar esta acción \\ \hline
		
		Poscondiciones
		& La información asociada al fondo autorizado para ejecutar pagos menores debe quedar almacenada en la base de datos \\ \hline
		
		Referencia
		& Paquete Configuración. RF 2.3 \\ \hline
		
		Prioridad
		& Crítico \\		
		\hline
	\end{supertabular}
	\caption[Descripción del caso de uso Modificar fondo autorizado para pagos menores]{Descripción del caso de uso Modificar fondo autorizado para pagos menores}
	\label{table:CU_ModFPM}
\end{table}

\begin{table}[H]
	\sf
	\begin{supertabular}{|p{4.5cm}|p{12cm}|}
		\hline
		\textbf{Nombre del caso de uso}
		& \textbf{Insertar saldos iniciales desde documento arqueo de caja} \\ \hline
		
		Actor
		& Jefe de Departamento \\ \hline
		
		Resumen
		& El caso de uso se inicia cuando el actor especificado desea fijar los saldos iniciales con que trabajará el sistema sobre la base de un arqueo de caja efectuado. El caso de uso termina cuando los saldos han sido introducidos satisfactoriamente \\ \hline
		
		Precondiciones
		& El usuario tiene que haberse autenticado en el sistema y contar con los privilegios necesarios \\ \hline
		
		Poscondiciones
		& La información asociada a los saldos iniciales del sistema debe quedar almacenada en la base de datos \\ \hline
		
		Referencia
		& Paquete Configuración. RF 2.5 \\ \hline
		
		Prioridad
		& Crítico \\		
		\hline
	\end{supertabular}
	\caption[Descripción del caso de uso Insertar saldos iniciales desde documento arqueo de caja]{Descripción del caso de uso Insertar saldos iniciales desde documento arqueo de caja}
	\label{table:CU_InsSI}
\end{table}

\begin{table}[H]
	\sf
	\begin{supertabular}{|p{4.5cm}|p{12cm}|}
		\hline
		\textbf{Nombre del caso de uso}
		& \textbf{Efectuar depósito de efectivo y cheques} \\ \hline
		
		Actor
		& Cajero \\ \hline
		
		Resumen
		& El caso de uso se inicia cuando el actor especificado desea realizar un depósito de efectivo y cheques al banco, una vez procesado los documentos procesados por estos conceptos (Recibo de Efectivo). El caso de uso termina cuando la información asociada a esta operación ha sido introducida satisfactoriamente \\ \hline
		
		Precondiciones
		& El usuario tiene que haberse autenticado en el sistema y contar con los privilegios necesarios. El usuario tiene que desglosar el efectivo introducido \\ \hline
		
		Poscondiciones
		& La información asociada al depósito de efectivo y cheques debe quedar almacenada en la base de datos. El reporte debe visualizarse en pantalla y ser exportable a PDF \\ \hline
		
		Referencia
		& Paquete Operaciones. RF 3.6 \\ \hline
		
		Prioridad
		& Crítico \\		
		\hline
	\end{supertabular}
	\caption[Descripción del caso de uso Efectuar depósito de efectivo y cheques]{Descripción del caso de uso Efectuar depósito de efectivo y cheques}
	\label{table:CU_EfectDepEf}
\end{table}

\begin{table}[H]
	\sf
	\begin{supertabular}{|p{4.5cm}|p{12cm}|}
		\hline
		\textbf{Nombre del caso de uso}
		& \textbf{Efectuar reembolso del fondo para dietas y pagos menores} \\ \hline
		
		Actor
		& Cajero \\ \hline
		
		Resumen
		& El caso de uso se inicia cuando el actor especificado desea realizar un reembolso del fondo efectivo usado para dietas y pagos menores, una vez seleccionados los documentos previamente procesados por estos conceptos (Vale para Pagos Menores y Anticipos a Justificar). El caso de uso termina cuando la información asociada a esta operación ha sido introducida satisfactoriamente \\ \hline
		
		Precondiciones
		& El usuario tiene que haberse autenticado en el sistema y contar con los privilegios necesarios. El sistema tiene que tener estos documentos sin haberse reembolsados para que la funcionalidad se ejecute \\ \hline
		
		Poscondiciones
		& La información asociada al Reembolso del Fondo para dietas y Pagos Menores debe quedar almacenada en la base de datos. El modelo debe visualizarse en pantalla y ser exportable a PDF \\ \hline
		
		Referencia
		& Paquete Operaciones. RF 3.4 \\ \hline
		
		Prioridad
		& Crítico \\		
		\hline
	\end{supertabular}
	\caption[Descripción del caso de uso Efectuar reembolso del fondo para dietas y pagos menores]{Descripción del caso de uso Efectuar reembolso del fondo para dietas y pagos menores}
	\label{table:CU_EfectReemb}
\end{table}

\begin{table}[H]
	\sf
	\begin{supertabular}{|p{4.5cm}|p{12cm}|}
		\hline
		\textbf{Nombre del caso de uso}
		& \textbf{Gestionar recibo de sellos del timbre} \\ \hline
		
		Actor
		& Cajero \\ \hline
		
		Resumen
		& El caso de uso se inicia cuando el actor especificado desea usar sellos para algún trámite. El cajero emite documento para la entrega de los sellos de acuerdo a la denominación y cantidad disponibles. El caso de uso termina cuando la información asociada a esta operación ha sido introducida satisfactoriamente \\ \hline
		
		Precondiciones
		& El usuario tiene que haberse autenticado en el sistema y contar con los privilegios necesarios. Los sellos deben introducirse en el sistema por denominación e importe \\ \hline
		
		Poscondiciones
		& La información asociada a la entrega de sellos debe quedar almacenada en la base de datos \\ \hline
		
		Referencia
		& Paquete Operaciones. RF 3.8 \\ \hline
		
		Prioridad
		& Secundario \\		
		\hline
	\end{supertabular}
	\caption[Descripción del caso de uso Gestionar recibo de sellos del timbre]{Descripción del caso de uso Gestionar recibo de sellos del timbre}
	\label{table:CU_GestSellos}
\end{table}

\begin{table}[H]
	\sf
	\begin{supertabular}{|p{4.5cm}|p{12cm}|}
		\hline
		\textbf{Nombre del caso de uso}
		& \textbf{Efectuar arqueo de caja} \\ \hline
		
		Actor
		& Jefe de Departamento \\ \hline
		
		Resumen
		& El caso de uso se inicia cuando el actor especificado desea efectuar un Arqueo de Caja al efectivo y documentos existentes, de acuerdo con la información que posee el sistema. El caso de uso termina cuando ha sido introducida correctamente la cantidad contada y los documentos verificados y esta no excede del fondo autorizado \\ \hline
		
		Precondiciones
		& El usuario tiene que haberse autenticado en el sistema y contar con los privilegios necesarios. Para realizar esta acción el sistema tiene que tener efectivo, cheques o sellos registrados y/o documentos de valor pendientes de ejecución. Introducir el efectivo físico. Realizar cálculo por denominación del efectivo existente \\ \hline
		
		Poscondiciones
		& Se emite el modelo de Arqueo de Caja una vez fijado el arqueo en pantalla, el modelo debe ser exportado a PDF \\ \hline
		
		Referencia
		& Paquete Operaciones. RF 3.5 \\ \hline
		
		Prioridad
		& Crítico \\		
		\hline
	\end{supertabular}
	\caption[Descripción del caso de uso Efectuar arqueo de caja]{Descripción del caso de uso Efectuar arqueo de caja}
	\label{table:CU_EfArqCaja}
\end{table}

\begin{table}[H]
	\sf
	\begin{supertabular}{|p{4.5cm}|p{12cm}|}
		\hline
		\textbf{Nombre del caso de uso}
		& \textbf{Modificar montos aprobados para gastos de viaje} \\ \hline
		
		Actor
		& Jefe de Departamento \\ \hline
		
		Resumen
		& El caso de uso se inicia cuando el actor especificado desea introducir los montos aprobados para gastos de viaje de acuerdo con la norma existente. El caso de uso termina cuando ha sido introducida correctamente los importes desglosados por tipo \\ \hline
		
		Precondiciones
		& El usuario tiene que haberse autenticado en el sistema y contar con los privilegios necesarios. El usuario debe contar con la información necesaria actualizada para acometer esta acción \\ \hline
		
		Poscondiciones
		& El sistema reflejará la información introducida la que quedará almacenada en la base de datos \\ \hline
		
		Referencia
		& Paquete Seguridad. RF 2.4 \\ \hline
		
		Prioridad
		& Secundario \\		
		\hline
	\end{supertabular}
	\caption[Descripción del caso de uso Modificar montos aprobados para gastos de viaje]{Descripción del caso de uso Modificar montos aprobados para gastos de viaje}
	\label{table:CU_ModMontos}
\end{table}

\begin{table}[H]
	\sf
	\begin{supertabular}{|p{4.5cm}|p{12cm}|}
		\hline
		\textbf{Nombre del caso de uso}
		& \textbf{Efectuar anticipo y liquidación de gastos de viajes} \\ \hline
		
		Actor
		& Cajero \\ \hline
		
		Resumen
		& El caso de uso se inicia cuando el actor especificado desea anticipar gastos de viaje o cuando quiere liquidar un anticipo en el sistema. El caso de uso termina cuando ha sido introducida la información necesaria, desglosada por monedas \\ \hline
		
		Precondiciones
		& El usuario tiene que haberse autenticado en el sistema y contar con los privilegios necesarios. El sistema calcula el efectivo a entregar por día de viaje \\ \hline
		
		Poscondiciones
		& El reporte debe visualizarse en pantalla y ser exportable a PDF \\ \hline
		
		Referencia
		& Paquete Operaciones. RF 3.3 \\ \hline
		
		Prioridad
		& Crítico \\		
		\hline
	\end{supertabular}
	\caption[Descripción del caso de uso Efectuar anticipo y liquidación de gastos de viajes]{Descripción del caso de uso Efectuar anticipo y liquidación de gastos de viajes}
	\label{table:CU_EfectAntic}
\end{table}

\begin{table}[H]
	\sf
	\begin{supertabular}{|p{4.5cm}|p{12cm}|}
		\hline
		\textbf{Nombre del caso de uso}
		& \textbf{Visualizar reportes} \\ \hline
		
		Actor
		& Usuario \\ \hline
		
		Resumen
		& El caso de uso se inicia cuando el actor especificado desea visualizar los reportes que emite el sistema, una vez que filtra el que desea. El caso de uso termina cuando ha sido mostrada la información solicitada \\ \hline
		
		Precondiciones
		& El usuario tiene que haberse autenticado en el sistema y contar con los privilegios necesarios y haber filtrado el reporte que necesita \\ \hline
		
		Poscondiciones
		& El reporte debe visualizarse en pantalla y ser exportable a PDF \\ \hline
		
		Referencia
		& Paquete Reportes. RF 4 \\ \hline
		
		Prioridad
		& Secundario \\		
		\hline
	\end{supertabular}
	\caption[Descripción del caso de uso Visualizar reportes]{Descripción del caso de uso Visualizar reportes}
	\label{table:CU_VisualRep}
\end{table}

\begin{table}[H]
	\sf
	\begin{supertabular}{|p{4.5cm}|p{12cm}|}
		\hline
		\textbf{Nombre del caso de uso}
		& \textbf{Efectuar Vales para Pagos Menores} \\ \hline
		
		Actor
		& Cajero \\ \hline
		
		Resumen
		& El caso de uso se inicia cuando el usuario desea efectuar un pago menor del efectivo existente en el sistema. El caso de uso termina cuando ha sido mostrado el reporte correspondiente \\ \hline
		
		Precondiciones
		& El usuario tiene que haberse autenticado en el sistema y contar con los privilegios necesarios. Debe desglosarse el efectivo pagado y no puede exceder el monto autorizado. Debe existir efectivo disponible para realizar el pago \\ \hline
		
		Poscondiciones
		& Se emitirá un modelo de Vale para Pagos Menores el cual se visualizará en pantalla y será exportable a PDF \\ \hline
		
		Referencia
		& Paquete Operaciones. RF 3.1 \\ \hline
		
		Prioridad
		& Crítico \\		
		\hline
	\end{supertabular}
	\caption[Descripción del caso de uso Efectuar Vales para Pagos Menores]{Descripción del caso de uso Efectuar Vales para Pagos Menores}
	\label{table:CU_ValePM}
\end{table}

\begin{table}[H]
	\sf
	\begin{supertabular}{|p{4.5cm}|p{12cm}|}
		\hline
		\textbf{Nombre del caso de uso}
		& \textbf{Efectuar Recibo de Efectivo} \\ \hline
		
		Actor
		& Cajero \\ \hline
		
		Resumen
		& El caso de uso se inicia cuando el actor especificado desea efectuar una entrega de efectivo del sistema del efectivo existente. El caso de uso termina cuando ha sido mostrado el reporte solicitado \\ \hline
		
		Precondiciones
		& El usuario tiene que haberse autenticado en el sistema y contar con los privilegios necesarios. Debe desglosarse el efectivo recibido \\ \hline
		
		Poscondiciones
		& Se emitirá un modelo de Recibo de Efectivo el cual se visualizará en pantalla y será exportable a PDF \\ \hline
		
		Referencia
		& Paquete Operaciones. RF 3.2 \\ \hline
		
		Prioridad
		& Crítico \\		
		\hline
	\end{supertabular}
	\caption[Descripción del caso de uso Efectuar Recibo de Efectivo]{Descripción del caso de uso Efectuar Recibo de Efectivo}
	\label{table:CU_EfectRecEfect}
\end{table}

\begin{table}[H]
	\sf
	\begin{supertabular}{|p{4.5cm}|p{12cm}|}
		\hline
		\textbf{Nombre del caso de uso}
		& \textbf{Visualizar trazas de usuarios} \\ \hline
		
		Actor
		& Administrador \\ \hline
		
		Resumen
		& El caso de uso se inicia cuando el actor especificado selecciona la opción trazas de usuario del sistema. El caso de uso finaliza con la visualización de las trazas de los usuarios \\ \hline
		
		Precondiciones
		& El usuario debe estar autenticado en el sistema y su rol debe ser administrador. Las trazas deben poder filtrarse por usuario y por fecha \\ \hline
		
		Poscondiciones
		& Se emitirá una ventana con la información solicitada \\ \hline
		
		Referencia
		& Paquete Seguridad. RF 1.8 \\ \hline
		
		Prioridad
		& Secundario \\		
		\hline
	\end{supertabular}
	\caption[Descripción del caso de uso Visualizar trazas de usuarios]{Descripción del caso de uso Visualizar trazas de usuarios}
	\label{table:CU_VisTrazaUs}
\end{table}

\begin{table}[H]
	\sf
	\begin{supertabular}{|p{4.5cm}|p{12cm}|}
		\hline
		\textbf{Nombre del caso de uso}
		& \textbf{Realizar salva de la base de datos} \\ \hline
		
		Actor
		& Administrador \\ \hline
		
		Resumen
		& El caso de uso se inicia cuando el actor especificado accede al sistema y selecciona la opción realizar la salva de la base de datos. El caso de uso finaliza con la visualización de un mensaje que notifica que la operación se realizó satisfactoriamente \\ \hline
		
		Precondiciones
		& El usuario debe estar autenticado en el sistema y su rol debe ser administrador.
		Debe estar definida la ruta de almacenamiento donde se va a guardar la salva la cual no puede estar en la partición del sistema operativo \\ \hline
		
		Poscondiciones
		& Debe quedar almacenada la salva de la información de la base de datos en la ruta previamente definida \\ \hline
		
		Referencia
		& Paquete Salvas. RF 5.2 \\ \hline
		
		Prioridad
		& Crítico \\		
		\hline
	\end{supertabular}
	\caption[Descripción del caso de uso Realizar salva de la base de datos]{Descripción del caso de uso Realizar salva de la base de datos}
	\label{table:CU_SalvaBD}
\end{table}

\begin{table}[H]
	\sf
	\begin{supertabular}{|p{4.5cm}|p{12cm}|}
		\hline
		\textbf{Nombre del caso de uso}
		& \textbf{Importar salva de la base de datos} \\ \hline
		
		Actor
		& Administrador \\ \hline
		
		Resumen
		& El caso de uso se inicia cuando el actor especificado desea restaurar una salva de la base de datos. El caso de uso finaliza con la visualización de un mensaje notificando que la operación se realizó de forma satisfactoria \\ \hline
		
		Precondiciones
		& El usuario debe estar autenticado en el sistema y su rol debe ser administrador.
		Debe existir al menos una salva de la base de datos en una ruta previamente definida, la cual debe ser seleccionada previamente \\ \hline
		
		Poscondiciones
		& La base de datos debe quedar actualizada con la información de la salva seleccionada \\ \hline
		
		Referencia
		& Paquete Salvas. RF 5.2 \\ \hline
		
		Prioridad
		& Auxiliar \\		
		\hline
	\end{supertabular}
	\caption[Descripción del caso de uso Importar salva de la base de datos]{Descripción del caso de uso Importar salva de la base de datos}
	\label{table:CU_ImportBD}
\end{table}

\begin{table}[H]
	\sf
	\begin{supertabular}{|p{4.5cm}|p{12cm}|}
		\hline
		\textbf{Nombre del caso de uso}
		& \textbf{Obtener ayuda sobre funcionalidades del sistema} \\ \hline
		
		Actor
		& Usuario \\ \hline
		
		Resumen
		& El caso de uso inicia cuando el actor especificado accede al sistema para visualizar la ayuda y lograr una mejor comprensión \\ \hline
		
		Precondiciones
		& El usuario debe estar autenticado en el sistema y contar con los permisos necesarios para efectuar esta operación \\ \hline
		
		Poscondiciones
		& La información mostrada debe coincidir con la seleccionada por el usuario \\ \hline
		
		Referencia
		& Paquete Reportes. RF 4.5 \\ \hline
		
		Prioridad
		& Auxiliar \\		
		\hline
	\end{supertabular}
	\caption[Descripción del caso de uso Obtener ayuda sobre funcionalidades del sistema]{Descripción del caso de uso Obtener ayuda sobre funcionalidades del sistema}
	\label{table:CU_ObtAyuda}
\end{table}

\begin{table}[H]
	\sf
	\begin{supertabular}{|p{4.5cm}|p{12cm}|}
		\hline
		\textbf{Nombre del caso de uso}
		& \textbf{Conceptos de documentos a procesar} \\ \hline
		
		Actor
		& Jefe de Departamento de Contabilidad \\ \hline
		
		Resumen
		& El caso de uso se inicia cuando el actor especificado desea gestionar los departamentos, la acción permite insertar, modificar y eliminar conceptos. Termina con un listado de los conceptos gestionados \\ \hline
		
		Precondiciones
		& El usuario debe estar autenticado en el sistema y contar con los permisos necesarios para efectuar esta operación \\ \hline
		
		Poscondiciones
		& La información mostrada debe coincidir con la seleccionada por el usuario \\ \hline
		
		Referencia
		& Paquete Configuración. RF 2.2 \\ \hline
		
		Prioridad
		& Secundario \\		
		\hline
	\end{supertabular}
	\caption[Descripción del caso de uso Conceptos de documentos a procesar]{Descripción del caso de uso Conceptos de documentos a procesar}
	\label{table:CU_ConcDoc}
\end{table}

\begin{table}[H]
	\sf
	\begin{supertabular}{|p{4.5cm}|p{12cm}|}
		\hline
		\textbf{Nombre del caso de uso}
		& \textbf{Control de anticipos a justificar} \\ \hline
		
		Actor
		& Especialista de Contabilidad \\ \hline
		
		Resumen
		& El caso de uso se inicia cuando el actor especificado desea confeccionar el estado de los anticipos a justificar. Termina con un reporte del Control de los Anticipos a Justificar \\ \hline
		
		Precondiciones
		& El usuario debe estar autenticado en el sistema y contar con los permisos necesarios para efectuar esta operación, además, deben existir anticipos de gastos de viajes liquidados o no para su visualización \\ \hline
		
		Poscondiciones
		& El reporte debe ser visualizado en pantalla y exportable a PDF \\ \hline
		
		Referencia
		& Paquete Reportes. RF 4.1 \\ \hline
		
		Prioridad
		& Auxiliar \\		
		\hline
	\end{supertabular}
	\caption[Descripción del caso de uso Control de anticipos a justificar]{Descripción del caso de uso Control de anticipos a justificar}
	\label{table:CU_ContAntic}
\end{table}

\begin{table}[H]
	\sf
	\begin{supertabular}{|p{4.5cm}|p{12cm}|}
		\hline
		\textbf{Nombre del caso de uso}
		& \textbf{Registrar cheque} \\ \hline
		
		Actor
		& Especialista de Contabilidad \\ \hline
		
		Resumen
		& El caso de uso inicia cuando se recibe un cheque que paga alguna factura de producto o servicio emitido por la entidad \\ \hline
		
		Precondiciones
		& El usuario debe estar autenticado en el sistema y contar con los permisos necesarios para efectuar esta operación, además, deben existir comprobantes de operaciones para exportar \\ \hline
		
		Poscondiciones
		& Debe aparecer el aviso de que los comprobantes fueron exportados \\ \hline
		
		Referencia
		& Paquete Operaciones. RF 3.7 \\ \hline
		
		Prioridad
		& Crítico \\		
		\hline
	\end{supertabular}
	\caption[Descripción del caso de uso Registrar cheque]{Descripción del caso de uso Registrar cheque}
	\label{table:CU_RegistCheque}
\end{table}

\begin{table}[H]
	\sf
	\begin{supertabular}{|p{4.5cm}|p{12cm}|}
		\hline
		\textbf{Nombre del caso de uso}
		& \textbf{Configurar parámetros de las salvas} \\ \hline
		
		Actor
		& Administrador \\ \hline
		
		Resumen
		& El caso de uso se inicia cuando el actor especificado desea configurar los parámetros de la salvan del sistema, para lo cual especifica el modo de ejecución y la ubicación donde se almacenarán las salvas, se tiene en cuenta que esta no puede estar en la partición del sistema operativo. El caso de uso finaliza con la visualización de la ventana principal de la aplicación \\ \hline
		
		Precondiciones
		& El usuario debe estar autenticado en el sistema y su rol debe ser administrador \\ \hline
		
		Poscondiciones
		& La ruta establecida no puede estar en el disco C de la PC donde se ejecutará la acción \\ \hline
		
		Referencia
		& Paquete Salvas. RF 5.3 \\ \hline
		
		Prioridad
		& Crítico \\		
		\hline
	\end{supertabular}
	\caption[Descripción del caso de uso Configurar parámetros de las salvas]{Descripción del caso de uso Configurar parámetros de las salvas}
	\label{table:CU_ConfigPSalva}
\end{table}

\section{Conclusiones parciales}
\begin{itemize}
	\item Se definieron como paquetes de la aplicación: seguridad, configuración,	operaciones y reportes
	\item Se obtuvieron 5 trabajadores del negocio que abarcan las funcionalidades	del sistema: “Administrador”, “Usuario”, “Cajero”, “Jefe de Departamento	de Contabilidad” y “Especialista en Contabilidad”
	\item Se clasificaron los casos de uso en “Crítico”, “Secundario” y “Auxiliar”,	para medir su impacto en el sistema
\end{itemize} % 

% % Para todos los capítulos, use la funcionalidad nueva chap{} en vez de chapter{}
% Esto hará que el texto en la izquierda superior de la página sea el mismo que el del capítulo

% Los Nombres de los capítulos se definen en Configuraciones/Administrativo.tex
\chap{\CapCuatro}

\sf
\section{Introducción}
\paragraph{}El diseño del \textit{software} es la última acción de la ingeniería correspondiente dentro de la actividad de modelado, la cual establece una plataforma para la construcción (generación de código y pruebas). Permite al ingeniero de \textit{software} modelar el sistema o producto que se va a construir, al convertir los requisitos del cliente en funcionalidades. Este flujo de trabajo garantiza que no se corran riesgos de construir un sistema inestable, propenso a fallos a la hora de realizar cambios, por muy pequeños que estos sean. Además, es la etapa donde se establece y fomenta la calidad del \textit{software} mediante el modelado de toda la solución \cite{goran_v._rational_2007}.

\section{Diagrama de clases de diseño}
\subsection{Diagrama de clases de diseño paquete configuraci\'{o}n}
Ver Anexos, figura \ref{fig:DCDPktConf}.

\subsection{Diagrama de clases de diseño paquete operaciones}
Ver Anexos, figura \ref{fig:DCDPktOper}.

\subsection{Diagrama de clases de diseño paquete reportes}
Ver Anexos, figura \ref{fig:DCDPktRep}.

\subsection{Diagrama de clases de diseño paquete seguridad}
Ver Anexos, figura \ref{fig:DCDPktSeg}.

\section{Estructuraci\'{o}n en capas}
\paragraph{}La arquitectura estructurada en capas permite que se encapsulen los datos de tal forma que se pueda tratar el código modularmente, sin correr el peligro de que, al modificar un código, deje de ser funcional la aplicación. Esto se logra al abstraer todo el código en diferentes capas, lo que permite además la interoperabilidad de las clases y la reutilización del código.
\paragraph{}El patrón de diseño Modelo Vista Controlador (MVC) es un estilo de arquitectura de \textit{software} que separa los datos de una aplicación, la interfaz de usuario, y la lógica de control en tres componentes distintos.
\paragraph{}Se trata de un modelo muy maduro y que ha demostrado su validez a lo largo de los años en todo tipo de aplicaciones, y sobre multitud de lenguajes y plataformas de desarrollo \cite{noauthor_modelo_2019}.

\begin{figure}[H] %Con el paquete float, se pone la opción [H] para que te salga la imagen donde la quieres
	\centering
	\includegraphics[width=15cm,height=8cm]{Figuras/MVC.png}
	\caption{Representación modelo de diseño arquitectónico Modelo Vista Controlador MCV}
	\label{fig:MVC}
\end{figure}

\paragraph{}A continuación se muestra el uso del patrón de diseño arquitectónico Modelo Vista Controlador en el sistema propuesto.

\begin{figure}[H] %Con el paquete float, se pone la opción [H] para que te salga la imagen donde la quieres
	\centering
	\includegraphics[width=17cm,height=8cm]{Figuras/MVCApp.png}
	\caption{Representación modelo de diseño arquitectónico Modelo Vista Controlador MCV}
	\label{fig:MVCApp}
\end{figure}

\paragraph{}Al combinar la arquitectura en capas con el enfoque MVC, se obtiene una arquitectura dividida en 3 capas, las cuales se tratarán según establece el patrón MVC. A continuación, en la en la Figura Estructuración en Capas se muestra el diagrama de estructuración en capas mediante el enfoque basado en Responsabilidad.
\paragraph{Capa Presentación:}Es la que presenta el sistema al usuario también conocida como interfaz de usuario, comunica la información y captura la información del usuario en un mínimo de proceso. Esta capa se comunica únicamente con la capa de negocio. 
\paragraph{Capa Lógica de Negocio:}En ella residen los programas que se ejecutan, se reciben las peticiones del usuario y se envían las respuestas tras el proceso, es aquí donde se establecen todas las reglas que deben cumplirse. Esta capa se comunica con la capa de presentación, para recibir las solicitudes y presentar los resultados, y con la capa de acceso a datos, para solicitar al gestor de base de datos almacenar o recuperar datos de él. 
\paragraph{Capa Acceso a Datos:}En ella residen los datos. Está formada por uno o más gestor de bases de datos que realiza todo el almacenamiento de datos, reciben solicitudes de almacenamiento o recuperación de información desde la capa de negocio.

\begin{figure}[H] %Con el paquete float, se pone la opción [H] para que te salga la imagen donde la quieres
	\centering
	\includegraphics[width=16cm,height=10cm]{Figuras/Estructura_capas.png}
	\caption{Estructuración en capas}
	\label{fig:Capas}
\end{figure}

\section{Patrones de diseño}
\paragraph{}Un patrón de diseño es una solución repetible a problemas típicos y recurrentes en el diseño del \textit{software}. Son soluciones basadas en la experiencia y que se ha demostrado que funcionan. No son un diseño terminado que puede traducirse directamente a código, sino más bien una descripción sobre cómo resolver el problema, la cual puede ser utilizada en diversas situaciones. Eric Gamma define los patrones de diseño como:
\paragraph{}“Un patrón de diseño es una descripción de clases y objetos comunicándose entre sí adaptada para resolver un problema de diseño general en un contexto particular” \cite{gamma_reutilizacion_2015}
\paragraph{}Por otra parte, Christopher Alexander describe a los patrones de conjunto con la Integridad Conceptual.
\paragraph{}"…  Cada patrón describe un problema que ocurre una y otra vez en nuestro entorno, para describir después el núcleo de la solución a ese problema, de tal manera que esa solución pueda ser usada más de un millón de veces sin hacerlo ni siquiera dos veces de la misma forma …" \cite{alexander_pattern_1997}.

\subsection{Fachada}
\paragraph{}El patrón FACADE simplifica el acceso a un conjunto de clases al proporcionar una única clase que todos utilizan para comunicarse con dicho conjunto de clases \cite{gamma_reutilizacion_2015}.
\subsubsection*{Ventajas}
\begin{itemize}
	\item Los clientes no necesitan conocer las clases que hay tras la clase FACADE.
	\item Se pueden cambiar las clases “ocultadas” sin necesidad de cambiar los clientes. Sólo hay que realizar los cambios necesarios en FACADE.
\end{itemize}

\paragraph{}En la aplicación se utiliza para facilitar el acceso a los datos, y la interacción con los mismos, como se muestra en la figura \ref{fig:Fachada}.

\begin{figure}%[H] %Con el paquete float, se pone la opción [H] para que te salga la imagen donde la quieres
	\centering
	\includegraphics[width=18cm,height=10cm]{Figuras/Fachada.png}
	\caption{Uso del patrón fachada}
	\label{fig:Fachada}
\end{figure}

\section{Principios de diseño}
\paragraph{}Para ayudar a desarrollar programas robustos, mantenibles y que se puedan modificar, existen varios principios de diseño que ayudan a desarrollar software de este tipo.
\paragraph{}Aunque la mayoría son aplicables a la programación orientada a objetos - el paradigma más extendido - algunos de ellos se pueden aplicar a otros tipos de programación \cite{rubenfa_doce_2014}.

\subsection{Interfaz de usuario}
\paragraph{}En todo desarrollo de \textit{software} la interfaz con la que se relaciona el usuario es de vital importancia para la aplicación, en este caso ha sido creada para aquel grupo de usuarios que tiene poca experiencia en el campo de la informática. El diseño desarrollado es en ventana, para seguir la misma estructura que la mayoría de las aplicaciones, es estable y uniforme. El diseño se rige fundamentalmente en color hueso, sin grandes pretensiones, con una imagen de fondo del Centro de Inmunología Molecular, en la misma aparece el nombre de la aplicación y su logo, en la parte inferior derecha se muestra el nombre del usuario que se encuentra al trabajar en el sistema, así como la fecha del día corriente. Se utilizan marcos rectangulares para los menús y submenús bien definidos, así como botones pequeños con una ayuda \textit{in situ} que permite al usuario ubicarse en la acción que tiene implícita cada uno. Para acceder al sistema el usuario debe ingresar su nombre de usuario y contraseña. Se utiliza una misma tipografía, forma y estilo en todas las ventanas, así como la simplicidad y consistencia de las mismas. 
\paragraph{}A continuación, se muestra la figura \ref{fig:AppMain}.

\begin{figure}[H] %Con el paquete float, se pone la opción [H] para que te salga la imagen donde la quieres
	\centering
	\includegraphics[width=17cm,height=10cm]{Figuras/APPMain.png}
	\caption{Interfaz gráfica de la aplicación}
	\label{fig:AppMain}
\end{figure}

\subsection{Formato de salida de los reportes}
\paragraph{}Todo sistema informático requiere salidas de la información que procesa para mostrar los resultados. En nuestro caso están encaminados a obtener información sobre las operaciones que procesa la caja y están basados en las normas emitidas por el Ministerio de Finanzas y Precios. Para la gestión de los reportes de la aplicación se usa \textit{IReport} debido a los beneficios que ofrece en cuanto a la generación de reportes y la integración con el lenguaje de programación que se eligió. El formato de salida de los reportes que se generan está en correspondencia con lo establecido en una norma contable\cite{noauthor_resolucion_2007}.
\paragraph{}A continuación, se muestra la estructura estándar que siguen los reportes de la aplicación:
\begin{itemize}
	\item Encabezado con nombre de la empresa.
	\item Encabezado con nombre del reporte.
	\item Fecha de emisión.
	\item Datos específicos del reporte.
\end{itemize}

\paragraph{}Para ilustrar lo anterior, se muestran ejemplos del formato de salida de los reportes en la figura \ref{fig:ReporteAnticipo} y la figura \ref{fig:ReporteControlAnticipo}.

\begin{figure}%[H] %Con el paquete float, se pone la opción [H] para que te salga la imagen donde la quieres
	\centering
	\includegraphics[width=17cm,height=10cm]{Figuras/ReporteAnticipo.png}
	\caption{Formato de salida del reporte Anticipo y Liquidación de Gastos de Viaje}
	\label{fig:ReporteAnticipo}
\end{figure}

\begin{figure}%[H] %Con el paquete float, se pone la opción [H] para que te salga la imagen donde la quieres
	\centering
	\includegraphics[width=17cm,height=4cm]{Figuras/ReporteControlAnticipos.png}
	\caption{Formato de salida del reporte Control de Anticipos a Justificar}
	\label{fig:ReporteControlAnticipo}
\end{figure}

\subsection{Principios de codificación}
\paragraph{}Camel Case es un principio de diseño perteneciente a la administración de configuración de cambios que establece un estándar para los nombres de las clases y los métodos y atributos de un sistema.
\paragraph{}Para la implementación de la aplicación se utilizó el principio \textit{Lower Camel Case} (figura \ref{fig:LowerCamelCase}) para los nombres de los atributos y operaciones, mientras que \textit{Upper Camel Case} (figura \ref{fig:UpperCamelCase}) fue utilizado al nombrar las clases de la aplicación.

\begin{figure}[H] %Con el paquete float, se pone la opción [H] para que te salga la imagen donde la quieres
	\centering
	\includegraphics[width=14cm,height=4cm]{Figuras/LowerCamelCase.png}
	\caption{Principio de diseño \textit{Lower Camel Case}}
	\label{fig:LowerCamelCase}
\end{figure}

\begin{figure}[H] %Con el paquete float, se pone la opción [H] para que te salga la imagen donde la quieres
	\centering
	\includegraphics[width=8cm,height=5cm]{Figuras/UpperCamelCase.png}
	\caption{Principio de diseño \textit{Upper Camel Case}}
	\label{fig:UpperCamelCase}
\end{figure}

%\paragraph{}A continuación, en la figura se muestra el uso del patrón de diseño arquitectónico Modelo Vista Controlador (MCV).
%
%\begin{figure}[H] %Con el paquete float, se pone la opción [H] para que te salga la imagen donde la quieres
%	\centering
%	\includegraphics[width=18cm,height=10cm]{Figuras/MVCApp.png}
%	\caption{Aplicación de MVC en el paquete reportes}
%	\label{fig:MVCApp}
%\end{figure}

\section{Ayuda}
\paragraph{}El sistema cuenta con una ayuda simple que permite al usuario navegar por sus opciones, al facilitar su uso.
\paragraph{}A continuación, se muestra una pantalla donde hace referencia a la ayuda del sistema.

\begin{figure}[H] %Con el paquete float, se pone la opción [H] para que te salga la imagen donde la quieres
	\centering
	\includegraphics[width=18cm,height=10cm]{Figuras/Ayuda.png}
	\caption{Pantalla principal ayuda}
	\label{fig:Ayuda}
\end{figure}

\section{Tratamiento de errores}
\begin{figure}[H] %Con el paquete float, se pone la opción [H] para que te salga la imagen donde la quieres
	\centering
	\includegraphics[width=7cm,height=3cm]{Figuras/ErrorAuth.png}
	\caption{Fallo en autenticación}
	\label{fig:ErrorAuth}
\end{figure}

\paragraph{}La información de los errores que se pueden cometer tanto en la introducción de datos como en la validación de los mismos garantiza que la información introducida en el sistema sea lo más fiable posible. A continuación, se muestran algunas pantallas donde se expone el tratamiento de algunas validaciones implementadas en la aplicación.
\paragraph{}La figura \ref{fig:ErrorBill} muestra cómo se implementa la validación para que el usuario no pueda utilizar una cantidad de billetes mayor a la existente en la base de datos.

\begin{figure}[H] %Con el paquete float, se pone la opción [H] para que te salga la imagen donde la quieres
	\centering
	\includegraphics[width=14cm,height=11cm]{Figuras/ErrorBill.png}
	\caption{Error en validación de cantidad de billetes}
	\label{fig:ErrorBill}
\end{figure}

\paragraph{}La figura \ref{fig:CancelDoc} muestra el mensaje del sistema cuando se pretende cancelar un documento, ya que esta acción es irreversible. Nótese el cambio de color de los documentos cancelados con respecto al resto y que no desaparece de la pantalla, para que sea posible llevar el récord de la numeración y el control documental.

\begin{figure} %Con el paquete float, se pone la opción [H] para que te salga la imagen donde la quieres
	\centering
	\includegraphics[width=14cm,height=11cm]{Figuras/CancelDoc.png}
	\caption{Confirmar cancelación de un documento}
	\label{fig:CancelDoc}
\end{figure}

\section{Diseño de la base de datos}
\subsection{Modelo lógico de datos}
\paragraph{}A continuación, se muestran las clases persistentes necesarias para implementar las funcionalidades de la aplicación propuesta. El diagrama de clases persistentes o modelo lógico de datos, sirve de punto de partida para el diseño de la base de datos del sistema.
\paragraph{}Se explica de una forma más detallada como, con la ayuda de Hibernate se realizan estas acciones:
\begin{itemize}
	\item Por cada entidad, que son las clases estereotipadas como “Entity”, se crea una tabla en la base de datos con el mismo nombre. La herramienta genera de forma automática la llave para cada tabla, identificada como id (tipo de dato auto numérico).
	\item Por cada relación de asociación de uno a uno (1:1) se generará en el extremo de la	relación donde se haya especificado la agregación una llave foránea que se	corresponderá con la llave del otro extremo.
	\item Por cada relación de asociación de uno a muchos (1:M) se genera en la que	corresponde al extremo mucho (M) un atributo que funciona como llave foránea que	coincide con la llave (campo id) del extremo uno (1).
	\item Por cada relación de asociación de muchos a muchos (M:M), se crea una tabla	intermedia que tendrá como campos los identificadores de las entidades	relacionadas, y se comportarán como llaves foráneas, mientras que su campo llave	lo ocupará un atributo id con el objetivo de evitar la complejidad de las llaves	compuestas.
	\item Por cada asociación de herencia, se genera por cada clase involucrada una tabla. En las tablas correspondientes a las entidades hijas se crea un atributo que se corresponde con la llave en la tabla generada a partir de la entidad padre
\end{itemize}

\begin{landscape}
\begin{figure}[H] %Con el paquete float, se pone la opción [H] para que te salga la imagen donde la quieres
	\centering
	\includegraphics[width=25cm,height=17cm]{Figuras/ModLogBD.png}
	\caption{Modelo lógico de datos}
	\label{fig:ModLogBD}
\end{figure}

\subsection{Modelo físico de datos}
\paragraph{}El modelo físico de los datos constituye la representación física del modelo de clases persistentes visto anteriormente. Contiene el conjunto de tablas que conforman la base de datos (ver Anexos, figura \ref{fig:ModFisBD}.

\begin{figure}[H] %Con el paquete float, se pone la opción [H] para que te salga la imagen donde la quieres
	\centering
	\includegraphics[width=25cm,height=14cm]{Figuras/ModFisBD.png}
	\caption{Modelo físico de datos}
	\label{fig:ModFisBD}
\end{figure}

\end{landscape}

\subsection{Diagrama de despliegue}
\paragraph{}El diagrama de despliegue es un modelo de objetos que describe la distribución física del sistema mediante nodos. Cada uno de ellos representa un recurso de cómputo dedicado a una función determinada. Este modelo da la correspondencia entre la arquitectura de software y la del sistema hardware \cite{sparxsystems_diagrama_nodate}.
\paragraph{}A continuación se presenta el diagrama de despliegue del sistema propuesto con los requerimientos mínimos de \textit{software} y \textit{hardware} para cada nodo.

\begin{figure}[H] %Con el paquete float, se pone la opción [H] para que te salga la imagen donde la quieres
	\centering
	\includegraphics[width=16cm,height=12cm]{Figuras/DDespliegue.png}
	\caption{Diagrama de despliegue}
	\label{fig:DDespliegue}
\end{figure}

\subsubsection*{Descripción de cada nodo}
\paragraph{Nodo PC Cliente:}Tiene como propósito realizar las peticiones al nodo Servidor de Base de Datos. Constituye una generalización de las diferentes estaciones de trabajo donde se va a ejecutar la aplicación. Es donde el usuario del sistema hará uso de los servicios ofrecidos por el mismo. 
\paragraph{Nodo Servidor de Base de datos:}Su propósito es almacenar la información procesada por la aplicación. 
\paragraph{Nodo Impresora:}Su propósito es imprimir las informaciones que brinda el sistema a través de los reportes.

\section{Conclusiones parciales}
\begin{itemize}
	\item Se aplicó el patrón de diseño “Fachada” al diseño de las clases
	\item La ayuda del sistema guía al usuario a través de todas las funcionalidades implementadas en la aplicación
\end{itemize} % 

% % Para todos los capítulos, use la funcionalidad nueva chap{} en vez de chapter{}
% Esto hará que el texto en la izquierda superior de la página sea el mismo que el del capítulo

% Los Nombres de los capítulos se definen en Configuraciones/Administrativo.tex
\chap{\CapCinco}

\sf
\section{Introducción}
\paragraph{}En el desarrollo de un proyecto de software es necesario practicar pruebas para descubrir errores en el funcionamiento del sistema. Las pruebas pueden clasificarse en: validación, funcionales, caja negra, caja blanca y automatizadas \cite{pittet_distintos_nodate}. La necesidad de determinar la factibilidad económica de un proyecto de software es importante para determinar el costo beneficio de su desarrollo. En este tipo de industria el cálculo se realiza sobre la base de la determinación del tiempo estimado para el desarrollo del proyecto, los costos asociados y los beneficios tangibles e intangibles. En este capítulo se diseñan los casos de prueba utilizados en el desarrollo del sistema, así como la planificación del proyecto a partir de los casos de uso.

\section{Tipos de pruebas ejecutadas y justificación de su selección}
\paragraph{}Las pruebas son un instrumento importante en el proceso de desarrollo de \textit{software}, debido a que evalúan la calidad del producto de \textit{software}, detectan defectos y los documentan, validan los requerimientos planteados y su funcionamiento. Las pruebas tienen el objetivo de detectar la mayor cantidad de errores posible y solucionarlos \cite{pittet_distintos_nodate}.

\section{Tipos de pruebas existentes}
Existen varios tipos de prueba a aplicar al \textit{software} dentro de las que se encuentran, según muestra la tabla \ref{table:TiposPruebas} \cite{quijano_que_2018}.

\begin{table}[ht]
	\begin{tabular}{ |p{3cm}|p{2cm}|p{4cm}|p{4cm}|p{2cm}| }
		\hline
		\textbf{Tipo de prueba}
		& \textbf{Practicada sobre}
		& \textbf{Técnicas}
		& \textbf{Resultados}
		& \textbf{Realizadas por} \\ \hline
		
		\textbf{Validación}
		& Software
		& Evaluaciones, inspecciones y tutoriales
		& Comprobar que lo que se ha especificado coincide con lo que el usuario quería.
		& Probador \\ \hline
		
		\textbf{Funcionales}
		& Software
		& Ejecución, revisión y retroalimentación de funcionalidades. Se realizan sobre el sistema en funcionamiento
		& Comprobar que se cumple con lo especificado a través de casos de uso
		& Probador \\ \hline
		
		\textbf{Caja negra}
		& Software
		& Valores de entradas y salidas
		& Se desconoce la implementación del código por el probador
		& Probador \\ \hline
		
		\textbf{Caja blanca}
		& Software
		& Verificación de  interacción entre componentes que implementan un caso de uso
		& Interacción interna entre los componentes del sistema. Ligado al código fuente
		& Probador \\ \hline
		
		\textbf{Automatizadas}
		& Software
		& Automatizan actividades comunes que no requieren inteligencia humana
		& Rapidez, confiabilidad, programable
		& Analista \\
		
		\hline
	\end{tabular}
	\caption[Tipos de pruebas existentes a aplicar al software]{Tipos de pruebas existentes a aplicar al software \cite{pittet_distintos_nodate}}
	\label{table:TiposPruebas}
\end{table}

\paragraph{}Para probar el \textit{software} se aplicó la variante de caja negra, ya que se lleva a cabo sobre la interfaz del \textit{software}, al proporcionar al sistema valores de entrada para medir el comportamiento esperado de las salidas, se mide si el resultado obtenido es el esperado \cite{pittet_distintos_nodate}.
\paragraph{}En este caso se realizaron pruebas durante una semana a tres de los casos de uso del sistema.

\section{Diseño de los casos de prueba}
Caso de uso de prueba introducir saldos iniciales

Tipo de Prueba Realizada: Caja Negra

Condiciones de ejecución: Los importes límites de pagos menores tienen que estar presentes

\begin{table}[H]
	\sf
	\begin{supertabular}{|p{4cm}|p{4cm}|p{4cm}|p{4cm}|}
		\hline
		
		\multicolumn{2}{l}{CU Introducir saldos iniciales}
		&  \hspace{2 mm}
		&  \hspace{2 mm}
		\\ \hline
		
		\textbf{Nombre del Escenario}
		& \textbf{Descripción del escenario de la prueba}
		& \textbf{Respuesta del sistema}
		& Flujo central \\ \hline
		
		EC 1.1 Introducir saldos satisfactoriamente
		& Los datos son cargados por la aplicación al inicio de la misma
		& 4. El sistema carga los datos del último arqueo de caja como saldos iniciales
		& 1. Abrir la aplicación
		
		2. Ir a la opción Introducir saldos iniciales
		
		3. Teclear los datos del último arqueo impreso \\ \hline
		
		EC 1.2 Fallo en la carga de datos
		& Los datos no se cargan en la aplicación debido a que su suma excede del valor asignado a la caja
		& 4. El sistema muestra un mensaje de error con el texto "El Importe desglosado excede el máximo permitido"
		& 1. Abrir la aplicación
		
		2. Ir a la opción Introducir saldos iniciales
		
		3. Teclear los datos del último arqueo
		 \\
				
		\hline
	\end{supertabular}
	\caption[Caso de uso de prueba introducir saldos iniciales]{Caso de uso de prueba introducir saldos iniciales}
	\label{table:CUP_IntrodSaldos}
\end{table}

Caso de uso de prueba Emitir Modelo de Vale para Pagos Menores

Tipo de Prueba Realizada: Caja Negra

Condiciones de ejecución: Tiene que estar identificado el usuario en la aplicación

\begin{table}[H]
	\sf
	\begin{supertabular}{|p{4cm}|p{4cm}|p{4cm}|p{4cm}|}
		\hline
		
		\multicolumn{2}{l}{CU Emitir Vales para Pagos Menores}
		&  \hspace{2 mm}
		&  \hspace{2 mm}
		\\ \hline
		
		\textbf{Nombre del Escenario}
		& \textbf{Descripción del escenario de la prueba}
		& \textbf{Respuesta del sistema}
		& Flujo central \\ \hline
		
EC 1.1 Emitir modelo de vale para pagos menores
		& Los datos del vale se muestran por la aplicación en un formulario
		& 4. El sistema muestra un formulario de vale para pagos menores
		
		6. El sistema muestra el vale activo
		& 1. Abrir la aplicación
		
		2. Ir al menú Operaciones
		
		3. Seleccionar Vale para pagos menores
		
		5. Introducir el monto a extraer
		 \\ \hline
		
		EC 1.2 Emitir modelo de vale para pagos menores
		& Los datos del vale se muestran por la aplicación en una planilla cuando se carga una cantidad de efectivo no disponible
		& 4. El sistema muestra un formulario de vale para pagos menores
		
		6. El sistema muestra un mensaje de error con la cantidad de la denominación existente
		& 1. Abrir la aplicación
		
		2. Ir al menú Reportes
		
		3. Seleccionar la opción Vale para pagos menores
		
		5. Se teclea por denominación e importe la cantidad objeto de pago superior a la existente en el sistema
		\\
		
		\hline
	\end{supertabular}
	\caption[Caso de uso de prueba emitir modelo de vale para pagos menores]{Caso de uso de prueba emitir modelo de vale para pagos menores}
	\label{table:CUP_ModeloVPM}
\end{table}

Caso de uso de prueba Emitir Recibo de Efectivo

Tipo de Prueba Realizada: Caja Negra

Condiciones de ejecución: Tiene que estar identificado el usuario en la aplicación

\begin{table}[H]
	\sf
	\begin{supertabular}{|p{4cm}|p{4cm}|p{4cm}|p{4cm}|}
		\hline
		
		\multicolumn{2}{l}{CU Emitir Recibo de Efectivo}
		&  \hspace{2 mm}
		&  \hspace{2 mm}
		\\ \hline
		
		\textbf{Nombre del Escenario}
		& \textbf{Descripción del escenario de la prueba}
		& \textbf{Respuesta del sistema}
		& Flujo central \\ \hline
		
		EC 1.1 Emitir modelo de Recibo de Efectivo
		& Los datos del recibo se muestran por la aplicación en un formulario
		& 5. El sistema muestra un modelo de Recibo de Efectivo
		
		6. El sistema muestra un Recibo de Efectivo activo
		
		& 1. Abrir la aplicación
		
		2. Ir al menú Operaciones
		
		3. Seleccionar la opción Recibo de Efectivo
		
		5. El usuario desglosa la cantidad que entregará
		\\ \hline
		
		EC 1.2 Imprimir modelo de Recibo de Efectivo
		& Los datos del recibo se muestran para su impresión
		& 4. El sistema muestra un modelo de Recibo de Efectivo listo para su llenado
		
		6. El sistema muestra la opción de imprimir el recibo
		& 1. Abrir la aplicación
		
		2. Ir al menú Operaciones
		
		3. Seleccionar la opción Recibo de Efectivo
		
		5. Teclear los datos del recibo
		\\ \hline
		
		EC 1.3 Fallo en la impresora
		& El recibo no se puede imprimir debido a que no hay una impresora conectada a la computadora
		& 6. El sistema muestra un modelo de Recibo de Efectivo listo para su impresión
		
		7. El sistema muestra la opción de imprimir el recibo
		& 1. Abrir la aplicación
		
		2. Ir al menú Reportes
		
		3. Seleccionar la opción Recibo de Efectivo
		
		5. Teclear los datos del recibo
		\\
		
		\hline
	\end{supertabular}
	\caption[Caso de uso de prueba Emitir Recibo de Efectivo]{Caso de uso de prueba Emitir Recibo de Efectivo}
	\label{table:CUP_ERE}
\end{table}

Caso de uso de prueba Visualizar Reportes

Tipo de Prueba Realizada: Caja Negra

Condiciones de ejecución: Debe haber al menos una impresora configurada

\begin{table}[H]
	\sf
	\begin{supertabular}{|p{4cm}|p{4cm}|p{4cm}|p{4cm}|}
		\hline
		
		\multicolumn{2}{l}{CU Visualizar reportes}
		&  \hspace{2 mm}
		&  \hspace{2 mm}
		\\ \hline
		
		\textbf{Nombre del Escenario}
		& \textbf{Descripción del escenario de la prueba}
		& \textbf{Respuesta del sistema}
		& Flujo central \\ \hline
		
		EC 1.1 Visualizar reportes
		& Los reportes no se muestran, ni se guardan por la aplicación ni se imprimen
		& 5. El sistema muestra un mensaje de error con el texto que aparece en la figura
		
		6. El menú se cierra sin obtenerse el reporte
		
		& 1. Abrir la aplicación
		
		2. Ir a la opción Reportes
		
		3. Escoger el reporte a imprimir
		
		5. Dar clic a botón aceptar
		\\
		
		\hline
	\end{supertabular}
	\caption[Caso de uso de prueba emitir Recibo de Efectivo]{Caso de uso de prueba emitir Recibo de Efectivo}
	\label{table:CUP_VisRep}
\end{table}

\begin{figure}[H] %Con el paquete float, se pone la opción [H] para que te salga la imagen donde la quieres
	\centering
	\includegraphics[width=15cm,height=3cm]{Figuras/EscenarioCasoPrueba.png}
	\caption{Escenario 1.1 Caso de prueba visualizar reportes}
	\label{fig:EscenarioCasoPrueba}
\end{figure}

\paragraph{}Las pruebas realizadas permitieron verificar que se cumplieron los objetivos propuestos al inicio de la implementación y corregir los errores. Se realizaron dos iteraciones de pruebas, la primera iteración produjo seis errores que se desglosan a continuación:
\begin{enumerate}
	\item Caso de uso de prueba: Introducir Saldos Iniciales
	\begin{enumerate}
		\item El sistema no muestra la información de todas las denominaciones y tiene errores en el cálculo del importe.
		\item El sistema no muestra el mensaje de error.
	\end{enumerate}
	\item Caso de uso de prueba: Emitir Vale para Pagos Menores
	\begin{enumerate}
		\item El formulario muestra errores
		\item El sistema no muestra correctamente la cantidad
	\end{enumerate}
	\item Caso de uso de prueba: Emitir Recibo de Efectivo
	\begin{enumerate}
		\item El sistema presenta errores en el formulario
		\item El sistema no muestra la opción de imprimir
	\end{enumerate}
\end{enumerate}

\paragraph{}Estos errores fueron corregidos mediante la modificación de código fuente en la aplicación.

\section{Esfuerzo del Implementador del sistema}
\begin{table}[ht]
	\begin{tabular}{ |p{4cm}|p{2.5cm}|p{1.5cm}|p{3cm}|p{2cm}| }
		\hline
		\textbf{Tareas Desarrolladas}
		& \textbf{Cantidad Trabajador}
		& \textbf{Salario / Hora}
		& \textbf{Tiempo en Horas Consumidas}
		& \textbf{Total} \\ \hline
		
		Captura de requisitos
		& 1 & 4.25 & 48 h (6 días) & $\$204.00$ \\ \hline
		
		Diseño de la solución propuesta
		& 1 & 4.25 & 200 h (25 días) & $\$850.00$ \\ \hline
		
		Propuesta al usuario
		& 1 & 4.25 & 48 h (6 días) & $\$204.00$ \\ \hline
		
		Selección y estudio de las herramientas
		& 1 & 4.25 & 192 h (24 días) & $\$816.00$ \\ \hline
		
		Implementación y documentación
		& 1 & 4.25 & 1920 h (240 días) & $\$8,160.00$ \\ \hline
		
		Implantación y capacitación
		& 1 & 4.25 & 40 h (5 días) & $\$170.00$ \\ \hline
		
		Pruebas
		& 1 & 4.25 & 80 h (10 días) & $\$340.00$ \\ \hline
		
		\multicolumn{3}{r}{\textbf{Total}} \vline
		& 2,528 horas & $\$10,744.00$ \\
		
		\hline
	\end{tabular}
	\caption[Esfuerzo del Implementador del sistema]{Esfuerzo del Implementador del sistema}
	\label{table:esfuerzoImpl}
\end{table}

\paragraph{}En total se invirtieron en el proyecto 2,528 horas que son equivalentes a 11 meses, tiempo en que se desarrolló el proyecto y la documentación.
\begin{center}
	{ \fboxsep 12pt
		\fcolorbox {black}{white}{
			\begin{minipage}[t]{12cm}
				Estimando que:
				
				Gastos Directos = $\$$ $10,744.00$
				
				Gastos Indirectos = $\$$ $3.00$ por hora (Estimación gastos indirectos consumo diario de electricidad, depreciación medios de cómputos, conexión a internet, papel y otros insumos)
				
				2,528 horas x 3 = $\$$ $843.00$
				
				Total Calculado = $\$$ $11,587.00$ $CUP$
			\end{minipage}
		} }
\end{center}

%\section{Planificación basada en casos de uso}
%\paragraph{}La planificación mediante el análisis de Puntos de Casos de Uso se trata de un método de estimación del tiempo de desarrollo de un proyecto mediante la asignación de “pesos” a un cierto número de factores que lo afectan[33].
%
%\subsection{Cálculo de los puntos de casos de uso sin ajustar.}
%\paragraph{}El primer paso para realizar la planificación es el cálculo de los puntos de los casos de uso sin ajustar, para ello se utilizará la siguiente fórmula:
%
%\begin{center}
%	{ \fboxsep 12pt
%		\fcolorbox {black}{white}{
%			\begin{minipage}[t]{4cm}
%				$PCU = FPA + PCU$
%			\end{minipage}
%		} }
%\end{center}
%\paragraph{}Donde:  
%\begin{itemize}
%	\item $PCU$: Puntos de Casos de Uso sin ajustar.
%	\item $FPA$: Factor de Peso de los Actores sin ajustar.
%	\item $FPCU$: Factor de Peso de los Casos de Uso sin ajustar.
%\end{itemize}
%
%\subsection{Factor de Peso de los Actores sin ajustar}
%\paragraph{}Este valor es resultante de un análisis de los actores presentes en el sistema y su complejidad, que se establece de acuerdo a los criterios que indica la tabla \ref{table:FactorPeso}.
%\paragraph{}El sistema cuenta con 5 actores complejos: Administrador, Usuario, Cajero, Jefe de Departamento y Especialista en Contabilidad; los que cumplen con la clasificación de complejo. De la información anterior se obtiene el Factor de peso x Actores ($FPA$).
%
%\begin{table}[ht]
%	\begin{tabular}{ |p{2.5cm}|p{8cm}|p{1.5cm}|p{1.5cm}|p{1.5cm}| }
%		\hline
%		\textbf{Tipo de actor}
%		& \textbf{Descripción}
%		& \textbf{Factor de peso}
%		& \textbf{Actores}
%		& \textbf{Total} \\ \hline
%		
%		\textbf{Simple}
%		& Otro sistema que interactúa con el sistema a desarrollar mediante una interfaz de aplicación
%		& 1 & 0 & 0 \\ \hline
%		
%		\textbf{Medio}
%		& Otro sistema que interactúa con el sistema a desarrollar mediante un protocolo o una interfaz basada en texto
%		& 2 & 0 & 0 \\ \hline
%		
%		\textbf{Complejo}
%		& Una persona que interactúa con el sistema a desarrollar mediante una interfaz gráfica
%		& 3 & 6 & 18 \\
%		
%		\hline
%	\end{tabular}
%	\caption[Factor de Peso de los actores sin ajustar]{Factor de Peso de los actores sin ajustar}
%	\label{table:FactorPeso}
%\end{table}
%
%\begin{center}
%	{ \fboxsep 12pt
%		\fcolorbox {black}{white}{
%			\begin{minipage}[t]{6cm}
%				$FPA = \sum(FactorDePeso * Actores)$ \\
%				$FPA = 3 * 6$ \\
%				$FPA = 18$
%			\end{minipage}
%		} }
%\end{center}
%
%\subsection{$FPCU$: Factor de Peso de los Casos de Uso sin ajustar} 
%\paragraph{}El valor del factor de peso de los casos de uso sin ajustar es el resultado del análisis de la cantidad de Casos de Uso existentes en el sistema y la complejidad que estos tienen.
%\paragraph{}Para determinar una transacción se tiene en cuenta la secuencia atómica de actividades, las cuales se realizan completamente o no se realiza ninguna.
%\paragraph{}En la tabla \ref{table:FactorPesoCU} se muestra la determinación del Factor de Peso de los Casos de Uso sin ajustar.
%
%\begin{table}[ht]
%	\begin{tabular}{ | l | l | l | l | l | }
%		\hline
%		\textbf{Tipo de caso de uso}
%		& \textbf{Descripción}
%		& \textbf{Peso}
%		& \textbf{\#CU}
%		& \textbf{\#CU*Peso} \\ \hline
%		
%		\textbf{Simple}
%		& El caso de uso contiene de 1 a 3 transacciones
%		& 5 & 26 & 130 \\ \hline
%		
%		\textbf{Medio}
%		& El caso de uso contiene de 4 a 7 transacciones
%		& 10 & 6 & 60 \\ \hline
%		
%		\textbf{Complejo}
%		& El caso de uso contiene de 8 o más transacciones
%		& 15 & 0 & 0 \\
%		
%		\hline
%	\end{tabular}
%	\caption[Factor de Peso de los Casos de Uso sin ajustar]{Factor de Peso de los Casos de Uso sin ajustar}
%	\label{table:FactorPesoCU}
%\end{table}
%
%\begin{center}
%	{ \fboxsep 12pt
%		\fcolorbox {black}{white}{
%			\begin{minipage}[t]{5cm}
%				$FCPU = \sum(\#CU * Peso)$ \\
%				$FCPU = 190$
%			\end{minipage}
%		} }
%\end{center}
%
%\begin{center}
%	{ \fboxsep 12pt
%		\fcolorbox {black}{white}{
%			\begin{minipage}[t]{7cm}
%				De lo anteriormente tenemos: \\
%				$FPA=18$ \\
%				y \\
%				$FCPU = 190$ \\
%				entonces: \\
%				$PCU= FPA + FPCU = 18 + 190 = 208$ \\
%				O sea: \\
%				$PCU = 208$
%			\end{minipage}
%		} }
%\end{center}
%
%\section{Cálculo de puntos de casos de uso ajustados}
%\paragraph{}Una vez que se tenga cálculo de los puntos de casos de uso sin ajustar ($PCU$) se procede a calcular los puntos de casos de uso ajustados ($PCUA$) cuyo resultado es el resultado es el producto del Factor de Complejidad Técnica ($FCT$) por el Factor Ambiente ($FA$), como se muestra a continuación:
%\begin{center}
%	{ \fboxsep 12pt
%		\fcolorbox {black}{white}{
%			\begin{minipage}[t]{5cm}
%				$FCUA = PCU * FCT * FA$
%			\end{minipage}
%		} }
%\end{center}
%\paragraph{Factor de Complejidad Técnica:}“Este valor se estima mediante la cuantificación del peso de un grupo de factores que determinan la complejidad técnica del software. A cada factor se le asigna un valor de 0 a 5 de acuerdo con la relevancia, donde 0 significa un aporte irrelevante y 5 un aporte muy importante” [33].
%
%\begin{table}[ht]
%	\begin{tabular}{ | l |p{4cm}| l | l |p{2cm}|p{6cm}| }
%		\hline
%		\textbf{Factor}
%		& \textbf{Descripción}
%		& \textbf{Peso}
%		& \textbf{Valor}
%		& \textbf{Peso*Valor}
%		& \textbf{Comentario} \\ \hline
%		
%		\textbf{T1}
%		& Sistema distribuido
%		& 2 & 3 & 63
%		& El sistema para desarrollar exige conexión de red, y un servidor de base de datos \\ \hline
%		
%		\textbf{T2}
%		& Objetivos de performance o tiempo de respuesta
%		& 1 & 3 & 3
%		& El sistema tiene un tiempo de respuesta normal \\ \hline
%		
%		\textbf{T3}
%		& Eficiencia del usuario final
%		& 1 & 2 & 2
%		& Existen escasas restricciones de eficiencias en el usuario final \\ \hline
%		
%		\textbf{T4}
%		& Procesamiento interno complejo
%		& 1 & 2 & 2
%		& No existe un procesamiento complejo \\ \hline
%		
%		\textbf{T5}
%		& El código debe ser reutilizable
%		& 1 & 5 & 5
%		& Se desea que el código sea lo más reusable posible por las magnitudes que puede alcanzar el software \\ \hline
%		
%		\textbf{T6}
%		& Facilidad de instalación
%		& 0.5 & 2 & 1
%		& No requiere grandes complejidades para su instalación \\ \hline
%		
%		\textbf{T7}
%		& Facilidad de uso
%		& 0.5 & 5 & 2.5
%		& El sistema presenta una interfaz fácil de usar y es sencilla \\ \hline
%		
%		\textbf{T8}
%		& Portabilidad
%		& 2 & 5 & 10
%		& El sistema es multiplataforma \\ \hline
%		
%		\textbf{T9}
%		& Facilidad de cambio
%		& 1 & 4 & 4
%		& El sistema está diseñado por capas por lo que los cambios se pueden realizar de forma sencilla \\ \hline
%		
%		\textbf{T10}
%		& Concurrencia
%		& 1 & 3 & 3
%		& Pueden acceder varios usuarios a la vez \\ \hline
%		
%		\textbf{T11}
%		& Incluye objetivos especiales de seguridad
%		& 1 & 3 & 3
%		& La seguridad se trata a partir de los privilegios establecidos, la autenticación y encriptación de la contraseña \\ \hline
%		
%		\textbf{T12}
%		& Provee acceso directo a terceras partes
%		& 1 & 0 & 0
%		& No provee acceso a terceras partes \\ \hline
%		
%		\textbf{T13}
%		& Se requieren facilidades especiales de entrenamiento
%		& 1 & 3 & 3
%		& El sistema es fácil de usar, pero el usuario requiere entrenamiento en su manejo \\		
%
%		\hline
%	\end{tabular}
%	\caption[Cálculo del Factor de Complejidad Técnica]{Cálculo del Factor de Complejidad Técnica}
%	\label{table:CalcFCT}
%\end{table}
%
%\paragraph{}Después de haber obtenido estos valores se procede a calcular el Factor de Complejidad Técnica ($FCT$).
%\begin{center}
%	{ \fboxsep 12pt
%		\fcolorbox {black}{white}{
%			\begin{minipage}[t]{7cm}
%				$FCT = 0.6 + 0.01 * \sum (Peso (i) * Valor (i))$\\
%				$FCT = 0.6 + 0.01 * (44.5) = 0.6 + 0.445$\\
%				$FCT = 1.045$
%			\end{minipage}
%		}
%	}
%\end{center}
%
%\subsubsection*{FA = Factor de Ambiente }
%\paragraph{}El Factor de Ambiente (FA) está relacionado con las habilidades y el entrenamiento del grupo de desarrollo del proyecto, estos factores tienen impacto en las estimaciones de tiempo. Este cálculo es similar al del Factor de complejidad técnica, se trata de un conjunto de factores que se cuantifican con valores de 0 a 5.
%\paragraph{}En la tabla \ref{table:DetermFA} se muestra el significado y el peso de cada uno de los factores y los valores asignados a cada uno de ellos.
%
%\begin{table}[ht]
%	\begin{tabular}{ | l |p{4.5cm}| l | l | l |p{5.5cm}| }
%		\hline
%		\textbf{Factor}
%		& \textbf{Descripción}
%		& \textbf{Peso}
%		& \textbf{Valor}
%		& \textbf{Peso*Valor}
%		& \textbf{Comentario} \\ \hline
%		
%		\textbf{E1}
%		& Familiaridad con el modelo de proyecto utilizado
%		& 1.5 & 5 & 7.5
%		& Existe familiaridad con el modelo de proyecto utilizado \\ \hline
%		
%		\textbf{E2}
%		& Experiencia en la aplicación
%		& 0.5 & 5 & 2.5
%		& Se tiene experiencia en la aplicación \\ \hline
%		
%		\textbf{E3}
%		& Experiencia en orientación a objetos
%		& 1 & 5 & 5
%		& Se tiene experiencia en la programación orientada a objetos \\ \hline
%		
%		\textbf{E4}
%		& Capacidad del analista líder
%		& 0.5 & 5 & 2.5
%		& El líder tiene la suficiente capacidad para llevar a cabo el desarrollo de la aplicación \\ \hline
%		
%		\textbf{E5}
%		& Motivación
%		& 1 & 5 & 5
%		& Existe alta motivación \\ \hline
%		
%		\textbf{E6}
%		& Estabilidad de los requerimientos
%		& 2 & 4 & 8
%		& Los requerimientos de la aplicación son estables \\ \hline
%		
%		\textbf{E7}
%		& Personal a tiempo parcial
%		& -1 & 2 & -2
%		& Tiempo compartido entre estudio y desarrollo de la aplicación \\ \hline
%		
%		\textbf{E8}
%		& Dificultad del lenguaje de programación
%		& -1 & 2 & -2
%		& Se empleará Java, el cual es un lenguaje de programación común \\
%		
%		\hline
%		\end{tabular}
%		\caption[Determinación del Factor de Ambiente]{Determinación del Factor de Ambiente}
%		\label{table:DetermFA}
%\end{table}
%
%\paragraph{}Con los resultados de la tabla anterior se procede a calcular el $FA$ (Factor de Ambiente) cuyo resultado se expresa en la siguiente ecuación:
%\begin{center}
%	{ \fboxsep 12pt
%		\fcolorbox {black}{white}{
%			\begin{minipage}[t]{9cm}
%				$FA = 1.4 – 0.03 * \sum (Peso (i) * Valor asignado (i)$ \\
%				$FA = 1.4 – 0.03 * (7.5 + 2.5 + 5 + 2.5 + 5 + 8 -2 -2)$ \\
%				$FA = 1.4 – 0.03 * (26.5)$ \\
%				$FA = 1.4 – 0.795$ \\
%				$FA = 0.605$
%				\end{minipage}
%		}
%	}
%\end{center}
%
%\section{Puntos de Casos de Uso Ajustados (PCUA)}
%\paragraph{}Por último y al tomar los valores calculados anteriormente se obtienen los Puntos de Casos de Uso Ajustados (PCUA).
%\begin{center}
%	{ \fboxsep 12pt
%		\fcolorbox {black}{white}{
%			\begin{minipage}[t]{5cm}
%				$PCUA = PCU * FCT * FA$ \\
%				$PCUA = 208 * 1.045 * 0.605$ \\
%				$PCUA = 131.503$
%			\end{minipage}
%		}
%	}
%\end{center}
%
%\section{Estimación del Esfuerzo}
%\paragraph{}Para el cálculo del esfuerzo en hombres/horas se utiliza la siguiente ecuación:
%\begin{center}
%	{ \fboxsep 12pt
%		\fcolorbox {black}{white}{
%			\begin{minipage}[t]{3cm}
%				$E = PCUA * FC$
%			\end{minipage}
%		}
%	}
%\end{center}
%
%Donde:
%\begin{itemize}
%	\item $E$ = Esfuerzo estimado en hombres/horas.
%	\item $PCUA$ = Puntos de Casos de Uso Ajustados.
%	\item $FC$ = Factor de Conversión.
%\end{itemize}
%
%\paragraph{}Para calcular el Factor de Conversión (FC) se analiza el Factor Ambiente y se observa cuáles de los que afectan el factor ambiente (E1…E6) se encuentran por debajo de la media (3) y que el resto (E7, E8) esté por encima de la media (3).
%\paragraph{}Para asignar la cantidad de horas/hombres se analiza que:
%\begin{itemize}
%	\item Si el total es dos o menos se utiliza el factor de conversión de veinte horas/hombre por Punto de Casos de Uso.
%	\item Si el total es tres o cuatro se utiliza como factor de conversión el valor de veinte y ocho horas/hombre por Punto de Casos de Uso.
%	\item Si el total es mayor o igual que cinco se recomienda efectuar cambios en el proyecto ya que existe un alto riesgo de fracaso.
%\end{itemize}
%
%\paragraph{}Al tomar como base lo expresado anteriormente se puede decir que:
%\begin{center}
%	{ \fboxsep 12pt
%		\fcolorbox {black}{white}{
%			\begin{minipage}[t]{8cm}
%				$FC = Horas/hombres * Punto de Casos de Uso$
%			\end{minipage}
%		}
%	}
%\end{center}
%\paragraph{}Al sustituir en la ecuación anterior obtenemos el siguiente resultado:
%\begin{center}
%	{ \fboxsep 12pt
%		\fcolorbox {black}{white}{
%			\begin{minipage}[t]{5cm}
%				$E = PCUA * FC$ \\
%				$E = 131.503 * 20 horas$ \\
%				$E = 2630.06 horas/hombres$
%			\end{minipage}
%		}
%	}
%\end{center}
%
%\paragraph{}La tabla \ref{table:EsfReq} muestra el cálculo del esfuerzo requerido, al tomar como base los porcentajes especificados:
%
%\begin{table}[ht]
%	\centering
%	\begin{tabular}{ | l | l | l | }
%		\hline
%		\textbf{Actividad}
%		& \textbf{Porcentaje}
%		& \textbf{Horas/Hombre} \\ \hline
%		
%		\textbf{Análisis}
%		& 10\% & 657.51 \\ \hline
%		
%		\textbf{Diseño}
%		& 20\% & 1315.03 \\ \hline
%		
%		\textbf{Implementación}
%		& 40\% & 2630.06 \\ \hline
%		
%		\textbf{Pruebas}
%		& 15\% & 986.27 \\ \hline
%		
%		\textbf{Sobrecarga (otras actividades)}
%		& 15\% & 986.27 \\ \hline
%		
%		\textbf{Total}
%		& 100\% & 6575.14 \\
%		
%		\hline
%		\end{tabular}
%		\caption[Esfuerzo requerido para el desarrollo del proyecto]{Esfuerzo requerido para el desarrollo del proyecto}
%		\label{table:EsfReq}
%\end{table}
%
%\paragraph{}Este esfuerzo es el que se requiere para la implementación del sistema. Si se tiene en cuenta que la fase de implementación representa el 40\% del esfuerzo total para desarrollar el software, entonces se puede calcular el esfuerzo total de la siguiente manera:
%\begin{center}
%	{ \fboxsep 12pt
%		\fcolorbox {black}{white}{
%			\begin{minipage}[t]{5cm}
%				$E (Total) = E / 0.4$ \\
%				$E (Total) = 2630.06 / 0.4$ \\
%				$E (Total) = 6575.14$
%			\end{minipage}
%		}
%	}
%\end{center}
%
%\section{Estimación del tiempo de desarrollo del proyecto}
%\begin{center}
%	{ \fboxsep 12pt
%		\fcolorbox {black}{white}{
%			\begin{minipage}[t]{5cm}
%				$TDES (total) =$ $E (total) \over CH (total) $
%			\end{minipage}
%		}
%	}
%\end{center}
%Donde: 
%\begin{itemize}
%	\item $TDES$ = Tiempo de desarrollo
%	\item $CH$ = Cantidad de Hombres = 1
%\end{itemize}
%
%\begin{center}
%	{ \fboxsep 12pt
%		\fcolorbox {black}{white}{
%			\begin{minipage}[t]{6cm}
%				$TDES (total) =$ $E (total) \over CH (total)$ \\
%				$TDES (total) =$ $6575 horas/hombre \over 1 hombre$ \\
%				$TDES (total) = 6575 horas$
%			\end{minipage}
%		}
%	}
%\end{center}
%
%\paragraph{}Si consideramos que el mes tiene 190.6 horas laborales (8 horas de trabajo por 24 días laborables como promedio)
%\begin{center}
%	{ \fboxsep 12pt
%		\fcolorbox {black}{white}{
%			\begin{minipage}[t]{8cm}
%				$$TDES (total) = \frac{6575 horas}{190.6 horas por mes}$$ \\
%				$TDES (total) = 34.5 meses \approx 35 meses$ \\
%				$TDES (total) = \frac{35 meses}{12 meses (1 ano)} = 2.916 anos$
%			\end{minipage}
%		}
%	}
%\end{center}
%\paragraph{}Lo que significa que el tiempo de desarrollo del proyecto es aproximadamente de 2 años y 9 meses.
%
%\section{Estimación del costo de desarrollo del proyecto}
%\begin{table}[ht]
%	\centering
%	\begin{tabular}{ |p{7cm}| l | l | }
%		\hline
%		\textbf{Variable}
%		& \textbf{Fórmula}
%		& \textbf{Valor} \\ \hline
%		
%		\textbf{C (total): costo total}
%		& $C(total)=C(p)+Otros gastos$
%		& 25 871.31 \\ \hline
%		
%		\textbf{C (p): costo del proyecto}
%		& $C (p)= E (total) * CHH$
%		& 25 871.31 \\ \hline
%		
%		\textbf{E (total): esfuerzo total}
%		& 
%		& 6575 \\ \hline
%		
%		\textbf{CHH: costo por hombre hora}
%		& $CHH= K * THP$
%		& 3.9348 \\ \hline
%		
%		\textbf{K: Coeficiente que tiene en cuenta los costos indirectos}
%		& 
%		& 1.50 \\ \hline
%		
%		\textbf{THP: Tarifa Horaria Promedio}
%		& $THP=SP/190.6$
%		& 2.6232 \\ \hline
%		
%		\textbf{SP: salario promedio}
%		& $SP=\sum _{n=i}^{n}\frac{Si}{n}$
%		& 500 \\\hline
%		
%		\textbf{Si: salario de trabajador}
%		& 
%		& 1200 \\
%		
%		\hline
%	\end{tabular}
%	\caption[Estimación total del costo del proyecto]{Estimación total del costo del proyecto}
%	\label{table:EstTotProy}
%\end{table}
%
%\begin{table}[ht]
%	\centering
%	\begin{tabular}{ | l | l | l | }
%		\hline
%		\textbf{Indicadores}
%		& \textbf{Valor}
%		& \textbf{Unidades} \\ \hline
%		
%		\textbf{Esfuerzo}
%		& 6575
%		& horas/hombre \\ \hline
%		
%		\textbf{Tiempo}
%		& 34.5
%		& meses \\ \hline
%		
%		\textbf{Desarrollo}
%		& 25 871.31
%		& MN \\
%		
%		\hline
%	\end{tabular}
%	\caption[Valor de los indicadores del proyecto]{Valor de los indicadores del proyecto}
%	\label{table:ValIProy}
%\end{table}
%
%\paragraph{}El costo estimado de desarrollo del proyecto está dado por la siguiente fórmula:
%\begin{center}
%	{ \fboxsep 12pt
%		\fcolorbox {black}{white}{
%			\begin{minipage}[t]{5cm}
%				$C (Total) = E (Total) * CHH$
%			\end{minipage}
%		}
%	}
%\end{center}
%\paragraph{}Donde:  
%\begin{itemize}
%	\item $C (Total)$: Costo total del proyecto.
%	\item $E (Total)$: Esfuerzo total.
%	\item $CHH$: Costo por hombre-hora.
%	\item $CHH = K * THP$
%\end{itemize}
%\paragraph{}Donde:  
%\begin{itemize}
%	\item $K$: Coeficiente que tiene en cuenta los costos indirectos (1,5 y 2,0).
%	\item $THP$: Tarifa Horaria Promedio. Se calcula al dividir el salario promedio de las personas que trabajan en el proyecto entre 190.6 horas laborables.
%\end{itemize}
%\begin{center}
%	{ \fboxsep 12pt
%		\fcolorbox {black}{white}{
%			\begin{minipage}[t]{8cm}
%				$Salario promedio = \sum \frac{Salario de los desarrolladores}{Cantidad de desarrolladores}$
%			\end{minipage}
%		}
%	}
%\end{center}
%\begin{center}
%	{ \fboxsep 12pt
%		\fcolorbox {black}{white}{
%			\begin{minipage}[t]{5cm}
%				$CHH = 1.5 * 2.623 = 3.9348$
%			\end{minipage}
%		}
%	}
%\end{center}
%\paragraph{}El proyecto cuenta con un desarrollador que tiene como salario 1200 CUP como promedio, por lo que:
%\begin{center}
%	{ \fboxsep 12pt
%		\fcolorbox {black}{white}{
%			\begin{minipage}[t]{4.5cm}
%				$THP = \frac{1200}{190.6} = 6.29590$
%			\end{minipage}
%		}
%	}
%\end{center}
%\begin{center}
%	{ \fboxsep 12pt
%		\fcolorbox {black}{white}{
%			\begin{minipage}[t]{5.5cm}
%				$C (Total) = E (Total) * K * CHH$ \\
%				$C (Total) = 6575 * 6.29590$ \\
%				$C (Total) = 41 395.54 CUP$
%			\end{minipage}
%		}
%	}
%\end{center}
%
%\paragraph{}Si utilizamos el tipo de cambio doméstico 25 CUP x 1.00 CUC y para poder comparar el costo de este proyecto en relación con los existentes en el mercado y que fueron analizados en el capítulo 1 de este documento, tenemos:
%\begin{center}
%	{ \fboxsep 12pt
%		\fcolorbox {black}{white}{
%			\begin{minipage}[t]{5cm}
%				$C (Total) = 1655.82 CUC$
%			\end{minipage}
%		}
%	}
%\end{center}

\section{Beneficios tangibles e intangibles}
\paragraph{}
Los beneficios tangibles son ventajas que se pueden medir en dólares que se acreditan a la organización mediante el uso del sistema de información los ejemplos de beneficios tangibles son un aumento en la velocidad del procesamiento, acceso de otra forma a la información inaccesible, acceso a la información en una forma más oportuna, ventaja por el poder de cálculo de la computadora y las disminuciones en el tiempo de empleado necesario para cumplir las tareas específicas. Aún hay otros, aunque la medición no siempre es fácil, actualmente los beneficios tangibles se pueden medir en términos de ahorros en dólares, recursos o tiempo \cite{clavo_beneficios_2013}.

\subsection{Beneficios tangibles}
\begin{itemize}
	\item El CIM cuenta con una herramienta cuyo costo de desarrollo asciende a $\$$ $411,587.00$ $CUP$ ($\$$ $1655.82$ $CUC$).
	\item Esta herramienta provee a la entidad la capacidad de contar con históricos de movimientos de operaciones de caja.
	\item Se elimina el gasto por concepto de modelaje para las operaciones de caja el cual era pagado en CUC.
	\item Se obtiene un \textit{software Open Source}, integrado a los procesos que se desarrollan en la caja, y se evita a la empresa el invertir en la compra de \textit{software} y la licencia.
\end{itemize}

\subsection{Beneficios intangibles}
\paragraph{}Los beneficios intangibles incluyen mejorar el proceso de toma de decisiones, incrementar la exactitud ser más competitivo con el servicio al cliente, mantener una buena imagen del negocio e incrementar la satisfacción del trabajo para los empleados eliminando las tareas tediosas. Cómo puede juzgar de la lista dada, los beneficios intangibles son sumamente importantes y pueden tener consecuencias transcendentales para el negocio conforme relaciona a las personas fuera y dentro de la organización \cite{clavo_beneficios_2013}.
\begin{itemize}
	\item Se evita la pérdida de información mediante una organización eficiente de estos en un sistema informático.
	\item Se homogeniza la información emitida en la caja con los documentos aprobados por el MFP y la misma puede ser consultada en cualquier momento.
	\item Se pueden realizar arqueos sorpresivos sobre la base de la información con que se cuenta en la aplicación, en ambas monedas.
\end{itemize}

\section{Análisis de costos y beneficios}
\paragraph{}Con la implantación del software propuesto la Dirección de Economía del Centro de Inmunología Molecular contará con una herramienta capaz de humanizar las operaciones de tesorería de caja de la empresa, así como permitirá controlar los datos que se manejan en ella. Si se toma el costo estimado del proyecto $\$$       $11,567.00$ $CUP$ ($\$$ $463.48$ $CUC$) y para su puesta en marcha no se incurrirán en gastos adicionales, esta opción constituye un ahorro de dinero. Además, con el uso de esta herramienta se mantienen actualizados los movimientos de efectivo existente en caja en tiempo real.

\section{Conclusiones parciales}
\begin{itemize}
	\item Los costos que se incurren en el desarrollo del sistema, ascienden a $\$$  $11,567.00$ $CUP$ ($\$$ $463.48$ $CUC$) en un tiempo estimado de 11 meses al trabajar una sola persona en el desarrollo del proyecto.
	\item Se realizaron cuatro casos de prueba y ocho escenarios en dos iteraciones.
\end{itemize} % 

%% Conclusiones ---------------------------------------------------
% Para que no le ponga delante a las Conclusiones la palabra Capítulo
\chapter*{Conclusiones}
%Adicionar a la tabla de contenidos la entrada Conclusiones
\addcontentsline{toc}{chapter}{Conclusiones}

\begin{itemize}
	\item Lorem ipsum dolor sit amet, consectetur adipiscing elit, sed do eiusmod tempor incididunt ut labore et dolore magna aliqua. Ut enim ad minim veniam, quis nostrud exercitation ullamco laboris nisi ut aliquip ex ea commodo consequat. Duis aute irure dolor in reprehenderit in voluptate velit esse cillum dolore eu fugiat nulla pariatur. Excepteur sint occaecat cupidatat non proident, sunt in culpa qui officia deserunt mollit anim id est laborum
\end{itemize}

%% Recomendaciones ---------------------------------------------------
% Para que no le ponga delante a las Recomendaciones la palabra Capítulo
\chapter*{Recomendaciones}

\addcontentsline{toc}{chapter}{Recomendaciones}

\begin{enumerate}
	\item Lorem ipsum dolor sit amet, consectetur adipiscing elit, sed do eiusmod tempor incididunt ut labore et dolore magna aliqua. Ut enim ad minim veniam, quis nostrud exercitation ullamco laboris nisi ut aliquip ex ea commodo consequat. Duis aute irure dolor in reprehenderit in voluptate velit esse cillum dolore eu fugiat nulla pariatur. Excepteur sint occaecat cupidatat non proident, sunt in culpa qui officia deserunt mollit anim id est laborum
\end{enumerate}


%% Referencias bibliográficas ---------------------------------------------------
\backmatter % Pista para decirle a LaTeX que los siguientes 'capítulos' son Referencias bibliográficas
\label{Referencias}
\lhead{\emph{Referencias Bibliográficas}}  % Cambiar encabezado izquierdo a "Referencias bibliográficas"
\addtotoc{Referencias Bibliográficas}
\renewcommand\bibname{Referencias Bibliográficas}
% Se pueden usar como estilos: plain, apalike, alpha, abbrv, unsrt
\bibliographystyle{apalike}  % Usar el estilo de BibTeX para el formato de la Bibliografía
% plain, apalike, alpha, abbrv, unsrt.

\begingroup
    \raggedright
    \sloppy
    \bibliography{Referencias}  % La información de las referencias bibliográficas se almacenan en el fichero llamado "Referencias.bib"
\endgroup 

\clearpage

%% -----------------------------------------------------

\clearpage
\pagestyle{fancy} % Los estilos de encabezados de página han estado "vacíos" todo este tiempo, 
                  % ahora se usan los encabezados "bonitos" como se definió anteriormente
\setstretch{1.5} % Establecer el espaciado a 1.5, 
                 % esto hace que las siguientes tablas sean más fáciles de leer
\lhead{\emph{Glosario de Términos}}  % Establecer el encabezado superior izquierdo a "Glosario de Términos"

% Para generar el glosario de términos, ir a Herramientas > Órdenes > PDFLaTeX, luego Herramientas > Órdenes > MakeIndex > y por último Herramientas > Órdenes > PDFLaTeX
\addtotoc{Glosario de Términos}

\cleardoublepage

\chapter*{Glosario de Términos}

\textbf{TÉRMINO:} Explicación del término. 

%% Anexos -----------------------------------------------------
\addtocontents{toc}{\vspace{2em}} % Añadir un espacio en los Contenidos, por estética
\lhead{\emph{Anexos}}  % Cambiar el encabezado izquierdo a "Anexos"
\appendix % Pista para decirle a LaTeX que los siguientes 'capítulos' son Anexos

\addtotoc{Anexos}
\chapter*{Anexos}

\section*{Anexo A}
\paragraph{}Lorem ipsum dolor sit amet, consectetur adipiscing elit, sed do eiusmod tempor incididunt ut labore et dolore magna aliqua. Ut enim ad minim veniam, quis nostrud exercitation ullamco laboris nisi ut aliquip ex ea commodo consequat. Duis aute irure dolor in reprehenderit in voluptate velit esse cillum dolore eu fugiat nulla pariatur. Excepteur sint occaecat cupidatat non proident, sunt in culpa qui officia deserunt mollit anim id est laborum.
	
 % Título del Anexo

\end{document}  % Fin
%% ----------------------------------------------------------------